\documentclass[a4paper,11pt]{article}

\usepackage{amsmath,amssymb,amsthm,mathtools}
\usepackage{bm}
\usepackage{geometry}
\usepackage{enumitem}
\usepackage{hyperref}

\geometry{
    top=25mm,
    bottom=25mm,
    left=25mm,
    right=25mm
}

%============================================================
% Theorem environments
%============================================================

\newtheorem{definition}{定義}[section]
\newtheorem{theorem}[definition]{定理}
\newtheorem{proposition}[definition]{命題}
\newtheorem{lemma}[definition]{補題}
\newtheorem{corollary}[definition]{系}
\newtheorem{remark}[definition]{注記}

%============================================================
% Notation
%============================================================

\newcommand{\Solver}{\mathcal{S}}
\newcommand{\SubSolver}{\preceq}
\newcommand{\Type}{\operatorname{Type}}
\newcommand{\DType}{\mathrm{D}}
\newcommand{\CType}{\mathrm{C}}
\newcommand{\MType}{\mathrm{M}}
\newcommand{\Aset}{\mathcal{A}}
\newcommand{\Rset}{\mathcal{R}}
\newcommand{\Th}{\mathrm{Th}}
\newcommand{\N}{\mathbb{N}}
\newcommand{\R}{\mathbb{R}}
\newcommand{\Pset}{\mathcal{P}}

%============================================================
% Title
%============================================================

\title{ECA Solver 集合に関する一般解の構築 - (A) 〜 (H)・改訂版 v2}

\author{}
\date{}

\begin{document}

\maketitle

%################################################################
% (A)
%################################################################

\section{(A) Solver 集合の数学的構造}
\label{sec:A}

\subsection{(A.1) Solver 集合の定義}
\label{subsec:A1}

ECA における Solver 集合を

\begin{equation}
    \Solver
\end{equation}

とする。

各

\begin{equation}
    S\in\Solver
\end{equation}

は、対象を所定の公理選択、推論規則および構成制約のもとで
評価するための構造を表す。

各 Solver $S$ を

\begin{equation}
    \boxed{
        S=
        \left(
            \Aset_S,
            \Rset_S
        \right)
    }
    \label{eq:A-solver-structure}
\end{equation}

と表す。

ここで、

\begin{itemize}
    \item $\Aset_S$ は Solver $S$ において採用される公理選択の集合、
    \item $\Rset_S$ は推論規則および構成制約の集合
\end{itemize}

を表す。

したがって、

\begin{equation}
    \Solver
    =
    \left\{
        S
        \,\middle|\,
        S=
        \left(
            \Aset_S,
            \Rset_S
        \right)
    \right\}
    \label{eq:A-solver-set}
\end{equation}

とする。

ここで Solver は単なる数値集合ではなく、少なくとも
公理選択および推論・構成規則からなる数学的構造として扱う。

また、Solver によって表現される理論を与える公理写像を

\begin{equation}
    \mathcal{T}:
    \Solver
    \longrightarrow
    \Th
    \label{eq:A-theory-map}
\end{equation}

とする。

各 $S\in\Solver$ に対して

\begin{equation}
    \mathcal{T}(S)
\end{equation}

は、$\Aset_S$ および $\Rset_S$ によって表現される理論を表す。

したがって、

\begin{equation}
    \mathcal{T}(S)
    =
    \mathcal{T}
    \left(
        \Aset_S,\Rset_S
    \right)
    \label{eq:A-theory-decomposition}
\end{equation}

と表記できる。

ただし、矛盾その他の理由によって理論を定義できない場合には、
$\mathcal{T}(S)$ は未定義となり得る。

\begin{remark}
本節では、Solver 集合の集合論的・構造論的性質を扱う。
可測族、測度および評価空間の測度構造については本節では導入せず、
後続の節において独立に定義する。
\end{remark}


%------------------------------------------------------------
\subsection{(A.2) 部分 Solver の定義}
\label{subsec:A2}

Solver の内部に存在する構造を扱うため、
単なる部分集合と部分 Solver を区別する。

\begin{definition}[部分 Solver]
$S,S'\in\Solver$ に対して、

\begin{equation}
    S'\SubSolver S
\end{equation}

とは、$S'$ が $S$ の要素の一部から構成され、
かつ $S$ において定義された Solver の構造を必要な範囲で継承して
独立した Solver として扱えることをいう。
\end{definition}

特に、

\begin{equation}
    S'\SubSolver S
    \quad\Longrightarrow\quad
    S'\subseteq S
    \label{eq:A-subsolver-subset}
\end{equation}

を満たすものとする。

ただし、逆に

\begin{equation}
    S'\subseteq S
\end{equation}

であることだけから

\begin{equation}
    S'\SubSolver S
\end{equation}

が成立するとは限らない。

すなわち、部分 Solver は単なる部分集合ではなく、
Solver として必要な構造を備えた部分構造である。

\begin{remark}
この定義により、一つの Solver の内部に複数の Solver が存在する
という階層構造を許容する。

したがって、

\begin{equation}
    S_1\SubSolver S,
    \qquad
    S_2\SubSolver S
\end{equation}

であって、

\begin{equation}
    S_1\neq S_2
\end{equation}

となることを妨げない。
\end{remark}


%------------------------------------------------------------
\subsection{(A.3) Solver の要素に対する離散性・連続性}
\label{subsec:A3}

Solver $S\in\Solver$ の要素 $s\in S$ に対して、
その要素が持つ数学的構造に基づいて離散的性質および連続的性質を
定義する。

離散的性質を満たす要素の集合を

\begin{equation}
    D_S
    :=
    \left\{
        s\in S
        \,\middle|\,
        s\text{ が所定の離散的構造を満たす}
    \right\}
    \label{eq:A-discrete-elements}
\end{equation}

とする。

同様に、連続的性質を満たす要素の集合を

\begin{equation}
    C_S
    :=
    \left\{
        s\in S
        \,\middle|\,
        s\text{ が所定の連続的構造を満たす}
    \right\}
    \label{eq:A-continuous-elements}
\end{equation}

とする。

ここでいう「離散的」および「連続的」は、単なる集合の濃度だけから
定義するものではない。

例えば、集合の可算性・非可算性は

\begin{equation}
    |S|\leq\aleph_0
\end{equation}

または

\begin{equation}
    |S|>\aleph_0
\end{equation}

によって既存の集合論から判定できるが、
非可算性そのものを直ちに「連続性」と同一視してはならない。

したがって、Solver の離散性・連続性の判定には、
必要に応じて既存の数学における

\begin{itemize}
    \item 集合論、
    \item 位相論、
    \item 順序構造、
    \item 代数構造、
    \item 解析学
\end{itemize}

等の既存概念を用いるものとする。


%------------------------------------------------------------
\subsection{(A.4) 中間的性質の定義}
\label{subsec:A4}

Solver $S$ において、ある要素が離散的性質と連続的性質を
同時に持つ場合を考える。

\begin{definition}[中間的要素]
$s\in S$ が中間的であるとは、

\begin{equation}
    s\in D_S\cap C_S
    \label{eq:A-intermediate-element}
\end{equation}

が成立することをいう。
\end{definition}

したがって、Solver $S$ が中間的性質を持つとは、

\begin{equation}
    \boxed{
        D_S\cap C_S\neq\varnothing
    }
    \label{eq:A-intermediate-solver}
\end{equation}

が成立することと定義する。

これは、Solver が単に離散部分と連続部分を一つの集合に
含んでいるという意味ではない。

重要なのは、

\begin{equation}
    \exists s\in S
\end{equation}

が存在して、

\begin{equation}
    s\in D_S
    \quad\text{かつ}\quad
    s\in C_S
\end{equation}

を同時に満たすことである。


%------------------------------------------------------------
\subsection{(A.5) Solver の型集合}
\label{subsec:A5}

前節までの定義から、Solver $S$ に対して一意の三択を
強制する必要はない。

そこで、Solver $S$ が有する型を集合として記録する。

\begin{definition}[Solver の型集合]
Solver $S\in\Solver$ に対して、その型集合を

\begin{equation}
    \boxed{
        \Type(S)
        \subseteq
        \left\{
            \DType,\CType,\MType
        \right\}
    }
    \label{eq:A-type-set}
\end{equation}

とする。
\end{definition}

ここで、

\begin{align}
    \DType\in\Type(S)
    &\quad\Longleftrightarrow\quad
    S\text{ が離散的 Solver 構造を含む},
    \label{eq:A-type-discrete}
    \\[2mm]
    \CType\in\Type(S)
    &\quad\Longleftrightarrow\quad
    S\text{ が連続的 Solver 構造を含む},
    \label{eq:A-type-continuous}
    \\[2mm]
    \MType\in\Type(S)
    &\quad\Longleftrightarrow\quad
    D_S\cap C_S\neq\varnothing
    \quad\text{または、}
    \nonumber\\
    &\hspace{28mm}
    S\text{ が中間 Solver 構造を含む}.
    \label{eq:A-type-intermediate}
\end{align}

とする。

この定義により、

\begin{equation}
    \Type(S)=\{\DType\}
\end{equation}

だけでなく、

\begin{equation}
    \Type(S)=\{\DType,\MType\},
\end{equation}

あるいは

\begin{equation}
    \Type(S)
    =
    \{\DType,\CType,\MType\}
\end{equation}

のような場合も許容される。

したがって、Solver の内部構造によって
「ある部分では離散的、別の部分では中間的」
という状況も表現できる。


%------------------------------------------------------------
\subsection{(A.6) 部分 Solver による型判定}
\label{subsec:A6}

Solver の型を判定する際には、Solver 全体を一つの対象として
強制的に三分類するのではなく、その内部に存在する部分 Solver を
調べることができる。

$S'\SubSolver S$ である各部分 Solver $S'$ に対して、

\begin{equation}
    \Type(S')
    \subseteq
    \{\DType,\CType,\MType\}
\end{equation}

を判定する。

このとき、Solver $S$ の内部に現れる型の集合を

\begin{equation}
    \boxed{
        \Type_{\mathrm{sub}}(S)
        :=
        \left\{
            \tau
            \,\middle|\,
            \exists S'\SubSolver S,\;
            \tau\in\Type(S')
        \right\}
    }
    \label{eq:A-subsolver-type-set}
\end{equation}

と定義する。

したがって、

\begin{equation}
    \Type_{\mathrm{sub}}(S)
    =
    \{\DType,\MType\}
\end{equation}

であれば、$S$ は内部に少なくとも
離散的 Solver 構造と中間的 Solver 構造を含む。

また、

\begin{equation}
    \Type_{\mathrm{sub}}(S)
    =
    \{\DType,\CType,\MType\}
\end{equation}

であれば、$S$ の内部には三種類すべての Solver 構造が
存在する。

\begin{remark}
ここで $\Type(S)$ と $\Type_{\mathrm{sub}}(S)$ は区別する。

$\Type(S)$ は $S$ 自身が持つ構造上の型を表すのに対し、
$\Type_{\mathrm{sub}}(S)$ は $S$ に含まれる部分 Solver の構造まで
調べた結果を表す。
\end{remark}


%------------------------------------------------------------
\subsection{(A.7) 型判定における既存数学の優先}
\label{subsec:A7}

Solver の型を判定するための数学的性質について、
既存の数学によって厳密に定義・判定できる場合には、
ECA 独自の量を導入する前に既存数学を用いる。

例えば、Solver または部分 Solver $S'$ の濃度については、

\begin{equation}
    |S'|\leq\aleph_0
\end{equation}

または

\begin{equation}
    |S'|>\aleph_0
\end{equation}

という通常の濃度論によって判定する。

また、位相構造が与えられている場合には、

\begin{equation}
    (S',\tau_{S'})
\end{equation}

について離散空間、連結空間その他の既存の位相的性質を
通常の位相論によって判定する。

この原則を、

\begin{equation}
    \boxed{
    \text{既存数学による判定}
    \;\longrightarrow\;
    \text{必要な場合のみ ECA の量による判定}
    }
    \label{eq:A-existing-math-priority}
\end{equation}

として採用する。

したがって、ECA における量 $B,K,I,C$ 等は、
既存数学によってSolverの性質を十分に記述できない場合に
その構造を評価するために導入する。


%------------------------------------------------------------
\subsection{(A.8) Solver の非可算性}
\label{subsec:A8}

Solver 集合 $S$ の非可算性を、

\begin{equation}
    |S|>\aleph_0
    \label{eq:A-uncountable}
\end{equation}

として表す。

ただし、本性質を任意の Solver に対して最初から仮定しない。

むしろ、ある部分 Solver

\begin{equation}
    S'\SubSolver S
\end{equation}

が存在して、

\begin{equation}
    |S'|>\aleph_0
    \label{eq:A-uncountable-subsolver}
\end{equation}

が既存数学から証明できるならば、

\begin{equation}
    S'\subseteq S
\end{equation}

より、

\begin{equation}
    |S|
    \geq
    |S'|
    >
    \aleph_0
\end{equation}

が従う。

したがって、

\begin{proposition}[部分 Solver による非可算性の継承]
$S,S'\in\Solver$ に対して

\begin{equation}
    S'\SubSolver S
\end{equation}

かつ

\begin{equation}
    |S'|>\aleph_0
\end{equation}

ならば、

\begin{equation}
    \boxed{
        |S|>\aleph_0
    }
\end{equation}

である。
\end{proposition}

\begin{proof}
$S'\SubSolver S$ より、

\begin{equation}
    S'\subseteq S
\end{equation}

が成立する。

したがって、包含写像

\begin{equation}
    \iota:S'\hookrightarrow S
\end{equation}

は単射である。

よって濃度について

\begin{equation}
    |S'|\leq |S|
\end{equation}

が成立する。

仮定より

\begin{equation}
    |S'|>\aleph_0
\end{equation}

であるから、

\begin{equation}
    |S|
    \geq
    |S'|
    >
    \aleph_0
\end{equation}

となる。

したがって、

\begin{equation}
    |S|>\aleph_0
\end{equation}

が従う。
\end{proof}

\begin{remark}
この命題により、Solver 全体の非可算性を直接仮定する必要はなく、
非可算な部分 Solver の存在から全体の非可算性を導出できる。
\end{remark}


%------------------------------------------------------------
\subsection{(A.9) Solver の構造と型の同型不変性}
\label{subsec:A9}

Solver 間の構造を比較するため、Solver 同型を考える。

\begin{definition}[Solver 同型]
$S_1,S_2\in\Solver$ に対して、

\begin{equation}
    T:
    S_1
    \overset{\sim}{\longrightarrow}
    S_2
\end{equation}

が Solver 同型であるとは、$T$ が全単射であり、
Solver の構造を保存する写像であることをいう。

特に、公理選択および推論規則・構成制約に対応する構造が

\begin{equation}
    T_{\mathcal A}:
    \Aset_{S_1}
    \overset{\sim}{\longrightarrow}
    \Aset_{S_2},
\end{equation}

\begin{equation}
    T_{\mathcal R}:
    \Rset_{S_1}
    \overset{\sim}{\longrightarrow}
    \Rset_{S_2}
\end{equation}

として対応するものとする。
\end{definition}

Solver 同型が離散的・連続的構造を保存する場合、

\begin{equation}
    s\in D_{S_1}
    \quad\Longleftrightarrow\quad
    T(s)\in D_{S_2},
    \label{eq:A-discrete-isomorphism}
\end{equation}

および

\begin{equation}
    s\in C_{S_1}
    \quad\Longleftrightarrow\quad
    T(s)\in C_{S_2}
    \label{eq:A-continuous-isomorphism}
\end{equation}

が成立する。

このとき、

\begin{equation}
    D_{S_2}
    =
    T(D_{S_1}),
\end{equation}

\begin{equation}
    C_{S_2}
    =
    T(C_{S_1})
\end{equation}

である。

したがって、

\begin{equation}
    D_{S_1}\cap C_{S_1}
    \neq
    \varnothing
\end{equation}

ならば、

\begin{equation}
    D_{S_2}\cap C_{S_2}
    \neq
    \varnothing
\end{equation}

である。

よって中間的性質は Solver 同型によって保存される。


%------------------------------------------------------------
\subsection{(A.10) 型集合の同型不変性}
\label{subsec:A10}

前節の構造保存条件を満たす Solver 同型

\begin{equation}
    T:
    S_1
    \overset{\sim}{\longrightarrow}
    S_2
\end{equation}

に対して、

\begin{equation}
    \boxed{
        \Type(S_1)=\Type(S_2)
    }
    \label{eq:A-type-invariance}
\end{equation}

が成立することを Solver の型の同型不変性と呼ぶ。

同様に、部分 Solver の型についても、

\begin{equation}
    S_1'\SubSolver S_1
\end{equation}

に対して

\begin{equation}
    T(S_1')\SubSolver S_2
\end{equation}

が対応するならば、

\begin{equation}
    \Type(S_1')
    =
    \Type(T(S_1'))
\end{equation}

が成立する。

したがって、

\begin{equation}
    \Type_{\mathrm{sub}}(S_1)
    =
    \Type_{\mathrm{sub}}(S_2)
    \label{eq:A-subtype-invariance}
\end{equation}

も成立する。


%------------------------------------------------------------
\subsection{(A.11) 集合論的構造と測度論的構造の分離}
\label{subsec:A11}

本節で定義した構造は、Solver の集合論的・構造論的性質に限定する。

特に、

\begin{equation}
    S,\quad
    S'\SubSolver S,\quad
    D_S,\quad
    C_S,\quad
    \Type(S)
\end{equation}

および、それらに関する濃度・部分構造・同型性を
本節の対象とする。

一方、

\begin{equation}
    \Sigma,\qquad
    \mu,\qquad
    (S,\Sigma,\mu)
\end{equation}

のような可測構造および測度空間の構成は、
ここでは定義しない。

したがって、ECA の構造を

\begin{equation}
    \boxed{
    \text{集合論的・構造論的層}
    \longrightarrow
    \text{評価空間の構成}
    \longrightarrow
    \text{測度論的層}
    }
    \label{eq:A-layer-separation}
\end{equation}

として分離する。

この分離により、

\begin{equation}
    \text{Solver の集合論的性質}
\end{equation}

と

\begin{equation}
    \text{Solver に導入される測度}
\end{equation}

を同一の定義として混同することを避ける。


%------------------------------------------------------------
\subsection{(A.12) (B) 以降への接続}
\label{subsec:A12}

以上により、(A) では Solver 集合について、

\begin{equation}
    \boxed{
    \begin{array}{c}
    \text{Solver の基本構造}\\
    \Downarrow\\
    \text{部分 Solver}\\
    \Downarrow\\
    \text{要素の離散性・連続性}\\
    \Downarrow\\
    \text{中間的性質}\\
    \Downarrow\\
    \text{Solver の型集合}\\
    \Downarrow\\
    \text{既存数学による型判定}\\
    \Downarrow\\
    \text{必要に応じた非可算性の導出}\\
    \Downarrow\\
    \text{Solver 同型による構造保存}
    \end{array}
    }
    \label{eq:A-structure-summary}
\end{equation}

という構造を確定する。

この構造を基礎として、(B) では Solver の評価構造を受け取る
一般評価空間を構成する。

すなわち、

\begin{equation}
    \Solver
    \longrightarrow
    X
\end{equation}

という評価写像を導入し、その後の (C) および (D) において
可測構造および測度を定義する。

また、(E) では本節で定義した Solver 同型を用いて、
評価構造および一般解の Representation Independence を扱う。

\begin{equation}
    \boxed{
    \text{Solver 構造}
    \longrightarrow
    \text{評価空間}
    \longrightarrow
    \text{測度構造}
    \longrightarrow
    \text{Solver 同型による不変性}
    }
    \label{eq:A-to-E}
\end{equation}

という順序によって、(A) から (E) までの構造を接続する。


%------------------------------------------------------------
\subsection*{(A) における定義上の注意}

本節では、Solver の「離散・連続・中間」という性質を
単純な三分法として扱わない。

特に、

\begin{equation}
    \Type(S)
    \subseteq
    \{\DType,\CType,\MType\}
\end{equation}

とすることで、複数の性質を同一の Solver が持つ可能性を許容する。

また、Solver の内部に

\begin{equation}
    S_1\SubSolver S
\end{equation}

という部分 Solver が存在することを許容するため、
Solver の型を判定する際には、その Solver 自身だけでなく、
必要に応じて部分 Solver の構造を調べることができる。

したがって、Solver の型判定は、

\begin{equation}
    \boxed{
    S
    \longrightarrow
    \{S'\mid S'\SubSolver S\}
    \longrightarrow
    \text{既存数学による構造判定}
    \longrightarrow
    \Type(S)
    }
    \label{eq:A-type-determination}
\end{equation}

という手順で行う。

このとき、既存数学によって十分に判定可能な性質については
ECA 独自の式を導入せず、必要な場合にのみ ECA の評価量を用いる。

%################################################################
% (B)
%################################################################

\section{(B) 一般評価空間の構成}
\label{sec:B}

%------------------------------------------------------------
\subsection{(B.1) 評価構造の基本的考え方}
\label{subsec:B1}

(A) において Solver 集合

\begin{equation}
    \Solver
\end{equation}

および各 Solver

\begin{equation}
    S
    =
    \left(
        \Aset_S,
        \Rset_S
    \right)
\end{equation}

の構造を定義した。

また、Solver の内部構造に応じて、
離散的・連続的・中間的な性質を持つ部分 Solver が存在し得ることを
定義した。

(B) では、これらの Solver 構造から得られる評価情報を、
Solver 自体とは独立した評価空間上の構造として表現する。

したがって、まず評価空間を

\begin{equation}
    X
\end{equation}

とする。

ここで $X$ は単なる集合としてではなく、
Solver によって得られる評価構造を保持するための数学的空間として
扱う。

ただし、$X$ にどのような追加構造を与えるかは、
評価対象および評価方法に応じて定めるものとする。

したがって、本節では $X$ に対して必要以上の構造を
あらかじめ仮定しない。

\begin{remark}
(A) において集合論的・構造論的な Solver の定義を行ったのに対し、
(B) ではその Solver を評価するための空間を構成する。

したがって、

\begin{equation}
    \text{Solver の構造}
    \longrightarrow
    \text{評価構造}
    \longrightarrow
    \text{評価空間}
\end{equation}

という順序を採用する。

可測族や測度については、本節で必要となる場合を除いて
具体的に導入せず、測度論的性質については (C) および (D) において
扱う。
\end{remark}


%------------------------------------------------------------
\subsection{(B.2) 評価写像}
\label{subsec:B2}

Solver $S\in\Solver$ に対して、その Solver が有する評価構造を
評価空間 $X$ の元として表現する写像を考える。

これを

\begin{equation}
    \boxed{
        \Phi:
        \Solver
        \longrightarrow
        X
    }
    \label{eq:B-evaluation-map}
\end{equation}

とする。

すなわち、任意の $S\in\Solver$ に対して、

\begin{equation}
    \Phi(S)\in X
\end{equation}

である。

ここで重要なのは、$\Phi$ が Solver を直ちに
単一の実数へ写す写像ではないことである。

$\Phi$ はまず、

\begin{equation}
    S
    \longmapsto
    \Phi(S)
\end{equation}

によって Solver が持つ評価構造を $X$ 上に表現する。

したがって、最終的な数値評価が必要な場合には、
後述する別の評価写像を介してその値を取り出す。

\begin{equation}
    \Solver
    \xrightarrow{\Phi}
    X
    \longrightarrow
    \R
\end{equation}

という二段階の構造を基本とする。


%------------------------------------------------------------
\subsection{(B.3) 一般評価空間の構成条件}
\label{subsec:B3}

評価空間 $X$ は、Solver の評価構造を表現できることを要求する。

すなわち、評価に必要な情報を表す集合を

\begin{equation}
    \mathcal E(S)
\end{equation}

とする。

ここで $\mathcal E(S)$ は Solver $S$ から得られる
評価に必要なデータの集合または構造を表す。

このとき、評価空間 $X$ は少なくとも

\begin{equation}
    \Phi(S)
    \in X
\end{equation}

となるように構成され、

\begin{equation}
    \Phi(S)
    =
    \operatorname{EvalStruct}(S)
    \label{eq:B-evaluation-structure}
\end{equation}

と表現できる。

ここで

\begin{equation}
    \operatorname{EvalStruct}
\end{equation}

は Solver の構造から評価構造を取り出す写像を表す。

ただし、

\begin{equation}
    \operatorname{EvalStruct}
\end{equation}

の具体的な定義は、Solver の型および評価方法によって異なり得る。

したがって、(B) の一般定義では特定の座標系に限定しない。


%------------------------------------------------------------
\subsection{(B.4) ECA における評価量による具体的表現}
\label{subsec:B4}

ECA において評価構造が

\begin{equation}
    B,\qquad I,\qquad K
\end{equation}

等の量によって記述される場合、
評価空間をそれぞれの評価量を受ける空間の直積として
表現できる。

すなわち、

\begin{equation}
    X_B,\qquad
    X_I,\qquad
    X_K
\end{equation}

をそれぞれの評価量の値域とすると、

\begin{equation}
    \boxed{
        X
        =
        X_B
        \times
        X_I
        \times
        X_K
    }
    \label{eq:B-product-space}
\end{equation}

と構成できる場合がある。

この場合、評価写像は

\begin{equation}
    \boxed{
        \Phi(S)
        =
        \left(
            B(S),
            I_S(S),
            K(S)
        \right)
    }
    \label{eq:B-coordinate-evaluation}
\end{equation}

と表現できる。

ここで、$I_S(S)$ は評価量としての $I$ を表す。

これは、公理選択空間として用いられる記号 $I$ と混同しないための
記法である。

\begin{remark}
ECA の既存文書では、$I$ が

\begin{equation}
    I:\Solver\to[0,1]
\end{equation}

という評価量として用いられる一方、

\begin{equation}
    I\subseteq\mathcal P(S)
\end{equation}

という公理選択空間を表す記号としても用いられている。

本稿では両者を数学的に区別する必要があるため、
前者を評価量 $I_S(S)$、後者を公理選択空間 $I$ として扱う。

\end{remark}


%------------------------------------------------------------
\subsection{(B.5) Solver の型に応じた評価構造}
\label{subsec:B5}

(A) において定義した Solver の型集合

\begin{equation}
    \Type(S)
    \subseteq
    \{
        \DType,\CType,\MType
    \}
\end{equation}

に応じて、評価構造の表現形式が異なる場合がある。

まず、離散的評価構造を持つ部分 Solver

\begin{equation}
    S_D\SubSolver S
\end{equation}

について、その評価に必要な離散データを

\begin{equation}
    B_D(S)
\end{equation}

とする。

ECA の既存定式化では、公理選択 $i$ に対して

\begin{equation}
    B(i)\subseteq i\subseteq S
    \label{eq:B-discrete-subset}
\end{equation}

を満たす離散的要素集合 $B(i)$ が導入されている。

この構造を用いる場合、離散評価は

\begin{equation}
    E_{\mathrm d}(S)
    =
    \sum_{b\in B(i_S)}
    w_S(b)
    \label{eq:B-discrete-evaluation}
\end{equation}

という形で表現できる。

ここで $i_S$ は Solver $S$ に対応する公理選択であり、
$w_S$ は対応する重み関数である。

ただし、公理選択 $i_S$ の一意な復元および
$w_S$ の厳密な構成は (H) で扱う。

したがって、(B) ではこの式を
離散評価構造の具体例として位置付ける。


%------------------------------------------------------------
\subsection{(B.6) 連続的評価構造}
\label{subsec:B6}

連続的な評価構造を持つ部分 Solver

\begin{equation}
    S_C\SubSolver S
\end{equation}

については、適切な測度

\begin{equation}
    \mu_S
\end{equation}

を用いることで評価を表現できる場合がある。

ECA における既存の連続形式では、公理選択 $i_S$ に対して

\begin{equation}
    E_{\mathrm c}(S)
    =
    \mu_S(i_S)
    =
    \int_S
    \mathbf{1}_{i_S}(u)
    \,d\mu_S(u)
    \label{eq:B-continuous-evaluation}
\end{equation}

という形が用いられる。

ここで

\begin{equation}
    \mathbf{1}_{i_S}:S\to\{0,1\}
\end{equation}

は $i_S$ の指示関数である。

この式は、離散的な総和を連続的な積分によって表現する
ECA の評価構造に対応する。

ただし、

\begin{equation}
    \mu_S
\end{equation}

そのものの存在、可測性および正規化可能性については、
本節では仮定の範囲を超えて断定せず、(C) および (D) で扱う。


%------------------------------------------------------------
\subsection{(B.7) 中間 Solver に対する評価空間}
\label{subsec:B7}

中間 Solver とは、(A) において定義したように、
離散的性質と連続的性質を同時に持つ要素を含む Solver である。

すなわち、

\begin{equation}
    D_S\cap C_S
    \neq
    \varnothing
    \label{eq:B-intermediate-condition}
\end{equation}

が成立する。

この場合、評価空間は離散評価構造と連続評価構造の双方を
表現できなければならない。

したがって、必要に応じて

\begin{equation}
    X
    =
    X_{\mathrm d}
    \times
    X_{\mathrm c}
    \times
    X_{\mathrm{dc}}
    \label{eq:B-intermediate-space}
\end{equation}

のように構成することができる。

ここで、

\begin{itemize}
    \item $X_{\mathrm d}$ は離散評価構造を表す空間、
    \item $X_{\mathrm c}$ は連続評価構造を表す空間、
    \item $X_{\mathrm{dc}}$ は両構造の対応関係を表すために
          必要となる空間
\end{itemize}

である。

ただし、この直積表示も一般定義そのものではなく、
中間 Solver の評価構造が実際にこの分解を許す場合の具体的表現である。

重要なのは、

\begin{equation}
    \operatorname{Type}(S)
    =
    \{\DType,\CType,\MType\}
\end{equation}

のように複数の型を持つ Solver を排除せず、
その内部に存在する各評価構造を同一の評価空間上で表現可能にする
ことである。


%------------------------------------------------------------
\subsection{(B.8) 一般評価空間と一般解}
\label{subsec:B8}

評価空間 $X$ が構成されたとき、
Solver に対して評価構造を対応させる一般的な写像を

\begin{equation}
    \boxed{
        \mathcal G:
        \Solver
        \longrightarrow
        X
    }
    \label{eq:B-general-solution}
\end{equation}

とする。

ここで $\mathcal G$ は、Solver に最終的な数値を直接対応させる
写像ではない。

任意の $S\in\Solver$ に対して、

\begin{equation}
    \mathcal G(S)\in X
\end{equation}

であり、

\begin{equation}
    \mathcal G(S)
\end{equation}

は Solver $S$ が持つ評価構造を一般解として表現する。

したがって、

\begin{equation}
    \boxed{
        S
        \longmapsto
        \mathcal G(S)
        \in X
    }
    \label{eq:B-general-solution-structure}
\end{equation}

を一般解の基本構造とする。

\begin{remark}
ここで $\mathcal G$ の存在を無条件に主張しているわけではない。

(B) では一般解が存在する場合に満たすべき写像の型と、
その評価空間を定義する。

\begin{equation}
    \exists\,
    \mathcal G:
    \Solver\to X
\end{equation}

そのものの成立については、(D) において存在問題として扱う。
\end{remark}


%------------------------------------------------------------
\subsection{(B.9) 一般解と最終評価}
\label{subsec:B9}

評価空間 $X$ の元から最終的なスカラー評価を取り出す写像を

\begin{equation}
    \Psi:
    X
    \longrightarrow
    \R
    \label{eq:B-final-evaluation-map}
\end{equation}

とする。

例えば、

\begin{equation}
    X
    =
    X_B\times X_I\times X_K
\end{equation}

かつ

\begin{equation}
    x=(B,I,K)
\end{equation}

である場合に、

\begin{equation}
    \Psi(B,I,K)
    =
    BIK
    \label{eq:B-ECA-final-evaluation}
\end{equation}

と定義することができる。

このとき、

\begin{equation}
    g
    :=
    \Psi\circ\mathcal G
    \label{eq:B-composed-evaluation}
\end{equation}

と置けば、

\begin{equation}
    \boxed{
        g:
        \Solver
        \longrightarrow
        \R
    }
\end{equation}

であり、

\begin{equation}
    \boxed{
        g(S)
        =
        \Psi(\mathcal G(S))
    }
    \label{eq:B-scalar-evaluation}
\end{equation}

を得る。

したがって、評価構造全体は

\begin{equation}
    \boxed{
        \Solver
        \xrightarrow{\mathcal G}
        X
        \xrightarrow{\Psi}
        \R
    }
    \label{eq:B-evaluation-chain}
\end{equation}

という二段階の写像として表現される。


%------------------------------------------------------------
\subsection{(B.10) 評価写像と一般解の区別}
\label{subsec:B10}

評価写像

\begin{equation}
    \Phi:\Solver\to X
\end{equation}

と一般解

\begin{equation}
    \mathcal G:\Solver\to X
\end{equation}

は、同じ型を持つ場合があるが、その意味を同一視する必要はない。

$\Phi$ は Solver の評価構造を評価空間上に表現する写像であり、

\begin{equation}
    \Phi(S)
\end{equation}

は Solver $S$ に対して得られる評価構造を表す。

一方、

\begin{equation}
    \mathcal G(S)
\end{equation}

は、その評価構造を ECA における一般解として表現する。

両者が同一の構造を表す場合には、

\begin{equation}
    \mathcal G(S)
    =
    \Phi(S)
    \qquad
    (\forall S\in\Solver)
\end{equation}

として同一視できる。

しかし、一般解の定義として

\begin{equation}
    \mathcal G=\Phi
\end{equation}

を要求する必要はない。

したがって、本稿では両者を区別した上で、

\begin{equation}
    \Phi
    \quad\text{および}\quad
    \mathcal G
\end{equation}

が同じ評価構造を表現する場合には、その対応を別途明示する。


%------------------------------------------------------------
\subsection{(B.11) 任意の評価対象への一般化}
\label{subsec:B11}

評価対象そのものを特定の集合に固定しないため、
任意の評価対象を

\begin{equation}
    Y
\end{equation}

とする。

このとき、評価対象 $Y$ から評価空間 $X$ への写像を

\begin{equation}
    \Phi_Y:
    Y
    \longrightarrow
    X
    \label{eq:B-object-evaluation}
\end{equation}

として定義できる場合がある。

さらに、複数の評価対象を一つのクラスまたは集合

\begin{equation}
    \mathfrak Y
\end{equation}

として扱う場合には、

\begin{equation}
    \Phi:
    \mathfrak Y
    \longrightarrow
    X
    \label{eq:B-general-object-evaluation}
\end{equation}

という一般化された評価構造を考えることができる。

これにより、評価対象そのものを固定することなく、
共通の評価空間構造を用いて異なる対象を評価できる。

ただし、この一般化された写像が Solver の評価写像

\begin{equation}
    \Phi:\Solver\to X
\end{equation}

と同一のものであるとは限らない。

評価対象と Solver を同一視できる場合にのみ、
両者を同じ写像として扱う。


%------------------------------------------------------------
\subsection{(B.12) Representation Independence への接続}
\label{subsec:B12}

評価写像の具体的表現は、Solver の表現方法によって
異なる場合がある。

したがって、

\begin{equation}
    \Phi_1:
    \Solver_1
    \longrightarrow
    X_1,
    \qquad
    \Phi_2:
    \Solver_2
    \longrightarrow
    X_2
\end{equation}

が異なる写像であっても、
Solver 同型

\begin{equation}
    T:
    \Solver_1
    \overset{\sim}{\longrightarrow}
    \Solver_2
\end{equation}

によって対応する評価構造を比較できることが要求される。

評価空間が同一の場合には、

\begin{equation}
    \boxed{
        \Phi_2\circ T
        =
        \Phi_1
    }
    \label{eq:B-representation-independence}
\end{equation}

という可換性によって対応を記述できる。

評価空間そのものが異なる場合には、

\begin{equation}
    U:
    X_1
    \longrightarrow
    X_2
\end{equation}

という評価空間間の構造保存写像を導入し、

\begin{equation}
    \boxed{
        U\circ\Phi_1
        =
        \Phi_2\circ T
    }
    \label{eq:B-evaluation-space-correspondence}
\end{equation}

とする。

この性質が、後述する (E) における
Representation Independence の基礎となる。

ここで重要なのは、評価空間の同値性を単なる集合としての同型性に
限定しないことである。

必要なのは、Solver の評価構造を保存する対応である。


%------------------------------------------------------------
\subsection{(B.13) 可測構造への接続}
\label{subsec:B13}

評価空間 $X$ に対して測度論的構造が必要となる場合には、
適切な可測空間

\begin{equation}
    (X,\Sigma_X)
\end{equation}

を導入する。

このとき、

\begin{equation}
    \Phi:
    \Solver
    \longrightarrow
    (X,\Sigma_X)
\end{equation}

または

\begin{equation}
    \mathcal G:
    \Solver
    \longrightarrow
    (X,\Sigma_X)
\end{equation}

の可測性を検討することができる。

ただし、可測族

\begin{equation}
    \Sigma_X
\end{equation}

および測度

\begin{equation}
    \mu
\end{equation}

の具体的構成は、本節で無条件に仮定しない。

これらについては (C) および (D) において、
評価関数の可測性、非負性、正規化可能性および測度構造との
整合性を別途検討する。


%------------------------------------------------------------
\subsection{(B.14) (A) から (B) への構造的帰結}
\label{subsec:B14}

(A) で定義した Solver 構造から、本節では

\begin{equation}
    \boxed{
        S
        \longmapsto
        \text{評価構造}
        \longmapsto
        X
    }
\end{equation}

という構成を導入した。

特に、Solver の型によって

\begin{equation}
    \Type(S)
    \subseteq
    \{\DType,\CType,\MType\}
\end{equation}

が与えられる場合、

\begin{equation}
    \DType
    \longrightarrow
    \text{離散評価構造},
\end{equation}

\begin{equation}
    \CType
    \longrightarrow
    \text{連続評価構造},
\end{equation}

\begin{equation}
    \MType
    \longrightarrow
    \text{離散・連続双方を含む評価構造}
\end{equation}

という対応を考えることができる。

ただし、これらは型の存在から評価空間上の具体的な構造が
自動的に存在することを意味するものではない。

評価構造の具体的な存在および一意性については、
後続の (C) および (D)、さらに Solver 同型との関係については
(E) において検証する。


%------------------------------------------------------------
\subsection*{(B) における定義上の注意}

本節では、以下の構造を明確に区別する。

\begin{enumerate}
    \item Solver 集合
    \begin{equation}
        \Solver
    \end{equation}

    \item Solver から評価空間への評価写像
    \begin{equation}
        \Phi:\Solver\to X
    \end{equation}

    \item 評価構造を一般解として表現する写像
    \begin{equation}
        \mathcal G:\Solver\to X
    \end{equation}

    \item 評価空間から最終評価値を取り出す写像
    \begin{equation}
        \Psi:X\to\R
    \end{equation}

    \item 最終評価関数
    \begin{equation}
        g=\Psi\circ\mathcal G.
    \end{equation}
\end{enumerate}

したがって、

\begin{equation}
    \boxed{
        \Solver
        \xrightarrow{\mathcal G}
        X
        \xrightarrow{\Psi}
        \R
    }
\end{equation}

が一般解の基本的な評価構造となる。

また、ECA における具体的な評価式として

\begin{equation}
    \Psi(B,I,K)=BIK
\end{equation}

を採用できる場合には、

\begin{equation}
    g(S)
    =
    B(S)I_S(S)K(S)
\end{equation}

となる。

ただし、この具体式を一般評価空間そのものの定義と同一視せず、
ECA の評価量がこの座標系で表現可能な場合の具体的評価として扱う。

以上により、(A) で定義された Solver の集合論的・構造論的性質を
基礎として、(B) ではその評価構造を保持する一般評価空間を
構成する。

%################################################################
% (C)
%################################################################

\section{(C) 評価関数の測度論的性質}
\label{sec:C}

%------------------------------------------------------------
\subsection{(C.1) 本節の目的}
\label{subsec:C1}

(B) において、Solver 集合 $\Solver$ から一般評価空間 $X$ への
写像

\begin{equation}
    \mathcal G:
    \Solver
    \longrightarrow
    X
\end{equation}

および、評価空間から最終的な評価値を取り出す写像

\begin{equation}
    \Psi:
    X
    \longrightarrow
    \R
\end{equation}

を定義した。

したがって、Solver $S\in\Solver$ に対する最終評価関数を

\begin{equation}
    \boxed{
        g
        :=
        \Psi\circ\mathcal G
    }
    \label{eq:C-evaluation-function}
\end{equation}

と定義する。

すなわち、

\begin{equation}
    g:
    \Solver
    \longrightarrow
    \R
\end{equation}

であり、

\begin{equation}
    g(S)
    =
    \Psi\!\left(\mathcal G(S)\right)
\end{equation}

である。

(C) では、この評価関数 $g$ を離散和および積分によって
扱うために必要となる測度論的性質を定義する。

特に、

\begin{equation}
    \text{非負性},
    \qquad
    \text{可測性},
    \qquad
    \text{正規化可能性}
\end{equation}

をそれぞれ独立した条件として扱う。


%------------------------------------------------------------
\subsection{(C.2) 評価空間の可測構造}
\label{subsec:C2}

(B) において構成された評価空間 $X$ に対して、
その評価を測度論的に扱うための $\sigma$-代数を

\begin{equation}
    \Sigma_X
\end{equation}

とする。

このとき、

\begin{equation}
    \boxed{
        (X,\Sigma_X)
    }
    \label{eq:C-measurable-evaluation-space}
\end{equation}

を評価空間の可測構造と呼ぶ。

ここで $\Sigma_X$ は

\begin{equation}
    \Sigma_X\subseteq\mathcal P(X)
\end{equation}

を満たす $\sigma$-代数である。

すなわち、

\begin{align}
    &X\in\Sigma_X,\\
    &A\in\Sigma_X
        \Longrightarrow
        X\setminus A\in\Sigma_X,\\
    &A_n\in\Sigma_X\ (n\in\N)
        \Longrightarrow
        \bigcup_{n=1}^{\infty}A_n\in\Sigma_X
\end{align}

を満たす。

これにより、評価空間 $X$ における
「測度を与える対象となる集合」を明確に定めることができる。

\begin{remark}
ここで $\Sigma_X$ の定義と、$\Sigma_X$ 上の測度の定義を区別する。

すなわち、

\begin{equation}
    \Sigma_X
\end{equation}

は可測集合の族そのものであり、

\begin{equation}
    \mu_X:\Sigma_X\to[0,\infty]
\end{equation}

はその可測集合に値を与える測度である。

したがって、

\begin{equation}
    \text{可測構造}
    \neq
    \text{測度}
\end{equation}

である。
\end{remark}


%------------------------------------------------------------
\subsection{(C.3) Solver 集合上の可測構造}
\label{subsec:C3}

評価関数 $g$ を Solver 集合上の測度論的対象として扱うため、
Solver 集合 $\Solver$ に対して $\sigma$-代数

\begin{equation}
    \Sigma_{\Solver}
\end{equation}

を導入する。

これにより、

\begin{equation}
    \boxed{
        (\Solver,\Sigma_{\Solver})
    }
    \label{eq:C-solver-measurable-space}
\end{equation}

を Solver 集合の可測空間とする。

特に、(B) の一般解

\begin{equation}
    \mathcal G:
    \Solver\to X
\end{equation}

が可測写像となることを要求する場合、

\begin{equation}
    \mathcal G:
    (\Solver,\Sigma_{\Solver})
    \longrightarrow
    (X,\Sigma_X)
\end{equation}

について、任意の

\begin{equation}
    A\in\Sigma_X
\end{equation}

に対して

\begin{equation}
    \mathcal G^{-1}(A)
    \in
    \Sigma_{\Solver}
\end{equation}

が成立することを要求する。

すなわち、

\begin{equation}
    \boxed{
        \mathcal G
        \text{ は可測}
    }
    \label{eq:C-G-measurable}
\end{equation}

である。


%------------------------------------------------------------
\subsection{(C.4) 最終評価写像の可測性}
\label{subsec:C4}

評価空間から最終評価値への写像

\begin{equation}
    \Psi:
    X\to\R
\end{equation}

について、$\R$ に Borel $\sigma$-代数

\begin{equation}
    \mathcal B(\R)
\end{equation}

を与える。

このとき、

\begin{equation}
    \Psi:
    (X,\Sigma_X)
    \longrightarrow
    (\R,\mathcal B(\R))
\end{equation}

が可測であるとは、任意の

\begin{equation}
    A\in\mathcal B(\R)
\end{equation}

について

\begin{equation}
    \Psi^{-1}(A)
    \in
    \Sigma_X
\end{equation}

が成立することである。

したがって、

\begin{equation}
    \boxed{
        \Psi
        \text{ は可測}
    }
    \label{eq:C-Psi-measurable}
\end{equation}

とする。


%------------------------------------------------------------
\subsection{(C.5) 最終評価関数の可測性}
\label{subsec:C5}

(B) において

\begin{equation}
    g
    =
    \Psi\circ\mathcal G
\end{equation}

と定義した。

ここで $\mathcal G$ および $\Psi$ がともに可測であるならば、
合成写像 $g$ も可測である。

\begin{proposition}[評価関数の可測性]
\label{prop:C-measurability}

\begin{equation}
    \mathcal G:
    (\Solver,\Sigma_{\Solver})
    \longrightarrow
    (X,\Sigma_X)
\end{equation}

および

\begin{equation}
    \Psi:
    (X,\Sigma_X)
    \longrightarrow
    (\R,\mathcal B(\R))
\end{equation}

が可測であるとする。

このとき、

\begin{equation}
    g
    =
    \Psi\circ\mathcal G
\end{equation}

は

\begin{equation}
    g:
    (\Solver,\Sigma_{\Solver})
    \longrightarrow
    (\R,\mathcal B(\R))
\end{equation}

として可測である。
\end{proposition}

\begin{proof}

任意の Borel 集合

\begin{equation}
    A\in\mathcal B(\R)
\end{equation}

を取る。

$\Psi$ の可測性より、

\begin{equation}
    \Psi^{-1}(A)
    \in
    \Sigma_X
\end{equation}

である。

さらに $\mathcal G$ の可測性より、

\begin{equation}
    \mathcal G^{-1}
    \left(
        \Psi^{-1}(A)
    \right)
    \in
    \Sigma_{\Solver}
\end{equation}

である。

一方、

\begin{align}
    g^{-1}(A)
    &=
    (\Psi\circ\mathcal G)^{-1}(A)
    \\
    &=
    \mathcal G^{-1}
    \left(
        \Psi^{-1}(A)
    \right).
\end{align}

したがって、

\begin{equation}
    g^{-1}(A)
    \in
    \Sigma_{\Solver}.
\end{equation}

$A$ は任意の Borel 集合だったので、

\begin{equation}
    g
    \text{ は可測}
\end{equation}

である。

\end{proof}


%------------------------------------------------------------
\subsection{(C.6) 評価関数の非負性}
\label{subsec:C6}

評価関数 $g$ が非負であることを、

\begin{equation}
    \boxed{
        g(S)\geq0
        \qquad
        (\forall S\in\Solver)
    }
    \label{eq:C-nonnegative}
\end{equation}

によって定義する。

すなわち、

\begin{equation}
    g:
    \Solver
    \longrightarrow
    [0,\infty)
\end{equation}

である。

この条件は、$g$ を非負測度との積分および
正規化操作に用いるための基本条件となる。

特に、ECA の評価構造が

\begin{equation}
    \mathcal G(S)
    =
    \left(
        B(S),
        I_S(S),
        K(S)
    \right)
\end{equation}

であり、

\begin{equation}
    \Psi(B,I,K)
    =
    BIK
\end{equation}

である場合には、

\begin{equation}
    g(S)
    =
    B(S)I_S(S)K(S)
\end{equation}

となる。

したがって、

\begin{equation}
    B(S)\geq0,
    \qquad
    I_S(S)\geq0,
    \qquad
    K(S)\geq0
\end{equation}

が保証されるならば、

\begin{equation}
    g(S)\geq0
\end{equation}

が従う。


%------------------------------------------------------------
\subsection{(C.7) Solver 集合上の測度}
\label{subsec:C7}

Solver 集合上の測度論的評価を行うため、

\begin{equation}
    \mu_{\Solver}:
    \Sigma_{\Solver}
    \longrightarrow
    [0,\infty]
\end{equation}

を導入する。

$\mu_{\Solver}$ は、以下の条件を満たす測度とする。

\begin{align}
    &\mu_{\Solver}(\varnothing)=0,\\
    &\mu_{\Solver}(A)\geq0
        \qquad
        (A\in\Sigma_{\Solver}),\\
    &A_n\in\Sigma_{\Solver}
        \text{ かつ }
        A_n\cap A_m=\varnothing
        \ (n\neq m)
\end{align}

ならば、

\begin{equation}
    \mu_{\Solver}
    \left(
        \bigcup_{n=1}^{\infty}A_n
    \right)
    =
    \sum_{n=1}^{\infty}
    \mu_{\Solver}(A_n)
\end{equation}

が成立する。

ここで重要なのは、

\begin{equation}
    \Sigma_{\Solver}
\end{equation}

の存在と、

\begin{equation}
    \mu_{\Solver}
\end{equation}

の存在を別々に扱うことである。

したがって、可測構造を定義したことのみから、
特定の測度が一意に決定されるとは限らない。


%------------------------------------------------------------
\subsection{(C.8) 評価関数の積分}
\label{subsec:C8}

$g$ が

\begin{equation}
    (\Solver,\Sigma_{\Solver})
    \longrightarrow
    ([0,\infty),\mathcal B([0,\infty)))
\end{equation}

として可測であるとき、

\begin{equation}
    \int_{\Solver}
    g(S)\,d\mu_{\Solver}(S)
\end{equation}

を定義できる。

この積分を、

\begin{equation}
    \boxed{
        Z
        :=
        \int_{\Solver}
        g(S)\,d\mu_{\Solver}(S)
    }
    \label{eq:C-normalization-constant}
\end{equation}

と置く。

$Z$ は評価関数全体の測度論的総量を表す。


%------------------------------------------------------------
\subsection{(C.9) 正規化可能性}
\label{subsec:C9}

評価関数 $g$ が $\mu_{\Solver}$ に関して正規化可能であるとは、

\begin{equation}
    \boxed{
        0
        <
        \int_{\Solver}
        g(S)\,d\mu_{\Solver}(S)
        <
        \infty
    }
    \label{eq:C-normalizable}
\end{equation}

が成立することと定義する。

すなわち、

\begin{equation}
    0<Z<\infty
\end{equation}

である。

このとき、

\begin{equation}
    \boxed{
        \widehat g(S)
        :=
        \frac{g(S)}{Z}
    }
    \label{eq:C-normalized-evaluation}
\end{equation}

と定義する。


%------------------------------------------------------------
\subsection{(C.10) 正規化の命題}
\label{subsec:C10}

\begin{proposition}[評価関数の正規化]
\label{prop:C-normalization}

$g$ が非負可測関数であり、

\begin{equation}
    0
    <
    Z
    =
    \int_{\Solver}
    g(S)\,d\mu_{\Solver}(S)
    <
    \infty
\end{equation}

を満たすならば、

\begin{equation}
    \widehat g(S)
    :=
    \frac{g(S)}{Z}
\end{equation}

は

\begin{equation}
    \int_{\Solver}
    \widehat g(S)\,d\mu_{\Solver}(S)
    =
    1
\end{equation}

を満たす。
\end{proposition}

\begin{proof}

$0<Z<\infty$ より、

\begin{equation}
    \frac{1}{Z}
\end{equation}

は有限の正数である。

したがって、積分の線形性より、

\begin{align}
    \int_{\Solver}
    \widehat g(S)\,d\mu_{\Solver}(S)
    &=
    \int_{\Solver}
    \frac{g(S)}{Z}
    \,d\mu_{\Solver}(S)
    \\
    &=
    \frac{1}{Z}
    \int_{\Solver}
    g(S)\,d\mu_{\Solver}(S)
    \\
    &=
    \frac{Z}{Z}
    \\
    &=1.
\end{align}

よって、

\begin{equation}
    \int_{\Solver}
    \widehat g(S)\,d\mu_{\Solver}(S)
    =
    1
\end{equation}

が成立する。
\end{proof}


%------------------------------------------------------------
\subsection{(C.11) 離散評価との対応}
\label{subsec:C11}

Solver $S$ が有限または可算的な離散評価構造を持ち、
各要素 $a\in S$ に対して重み

\begin{equation}
    w(a)\geq0
\end{equation}

が与えられている場合には、

\begin{equation}
    \sum_{a\in S}w(a)
\end{equation}

によって評価を構成できる。

特に、公理選択 $i\subseteq S$ に対して

\begin{equation}
    B(i)\subseteq i\subseteq S
\end{equation}

である場合、ECA における離散的測度の具体形として

\begin{equation}
    \boxed{
        \mu(i)
        =
        \sum_{b\in B(i)}
        w(b)
    }
    \label{eq:C-discrete-measure}
\end{equation}

を採用できる。

このとき、

\begin{equation}
    w:S\to[0,\infty)
\end{equation}

が非負であり、

\begin{equation}
    \sum_{a\in S}w(a)
    =
    1
\end{equation}

を満たすならば、対応する重み付き評価は
正規化された離散評価となる。


%------------------------------------------------------------
\subsection{(C.12) 連続評価との対応}
\label{subsec:C12}

Solver またはその評価対象上に可測空間

\begin{equation}
    (S,\Sigma_S)
\end{equation}

と測度

\begin{equation}
    \mu_S:\Sigma_S\to[0,\infty]
\end{equation}

が与えられている場合を考える。

公理選択

\begin{equation}
    i\subseteq S
\end{equation}

が可測集合であるとき、その指示関数

\begin{equation}
    \mathbf 1_i:S\to\{0,1\}
\end{equation}

を用いて、

\begin{equation}
    \boxed{
        \mu_S(i)
        =
        \int_S
        \mathbf 1_i(s)\,d\mu_S(s)
    }
    \label{eq:C-continuous-measure}
\end{equation}

と表現できる。

これは離散的な総和による評価を、
連続的な測度による評価へ拡張したものである。

ただし、この式が成立するためには、

\begin{equation}
    i\in\Sigma_S
\end{equation}

すなわち $i$ が可測集合であることを要求する。


%------------------------------------------------------------
\subsection{(C.13) 測度そのものと測度空間の区別}
\label{subsec:C13}

評価構造において測度を対象とする場合と、
測度空間全体を対象とする場合を区別する。

まず、測度そのものを対象とする場合には、

\begin{equation}
    \mu:
    \Sigma\to[0,\infty]
\end{equation}

を評価対象とする。

一方、測度空間を対象とする場合には、

\begin{equation}
    (X,\Sigma_X,\mu_X)
\end{equation}

という三つ組全体を評価対象とする。

したがって、

\begin{equation}
    \mu
\end{equation}

と

\begin{equation}
    (X,\Sigma_X,\mu)
\end{equation}

は同一の対象として扱わない。

この区別は、後に連続 Solver における

\begin{equation}
    C_s^{(\mu)}
\end{equation}

と

\begin{equation}
    C_s^{(\mathrm{ms})}
\end{equation}

を区別する際にも必要となる。


%------------------------------------------------------------
\subsection{(C.14) ECA の評価式との接続}
\label{subsec:C14}

ECA において最終評価が

\begin{equation}
    \Psi(B,I,K)
    =
    BIK
\end{equation}

で与えられる場合には、

\begin{equation}
    g(S)
    =
    B(S)I_S(S)K(S)
\end{equation}

となる。

さらに、評価対象が連続的な変数 $x$ によって表される場合には、

\begin{equation}
    \boxed{
        T_{\mathrm{overlay}}
        =
        \int
        B(x)I(x)K(x)
        \,d\mu(x)
    }
    \label{eq:C-overlay-integral}
\end{equation}

という形を得る。

また、区間 $[0,1]$ 上の Lebesgue 測度を採用する場合には、

\begin{equation}
    \boxed{
        T_{\mathrm{overlay}}
        =
        \int_0^1
        B(x)I(x)K(x)
        \,dx
    }
    \label{eq:C-lebesgue-integral}
\end{equation}

となる。

ただし、これらの具体的積分表示は、
対応する関数の可測性および積分可能性が保証される場合にのみ
成立する。


%------------------------------------------------------------
\subsection{(C.15) (A)--(B) から (C) への構造的帰結}
\label{subsec:C15}

(A) において Solver 集合の集合論的・構造論的基礎を定め、
(B) において一般評価空間

\begin{equation}
    X
\end{equation}

および評価構造

\begin{equation}
    \Solver
    \xrightarrow{\mathcal G}
    X
    \xrightarrow{\Psi}
    \R
\end{equation}

を定義した。

(C) では、この構造に対して

\begin{equation}
    \Sigma_{\Solver},
    \qquad
    \Sigma_X,
    \qquad
    \mu_{\Solver},
    \qquad
    \mu_X
\end{equation}

等の測度論的構造を必要に応じて導入し、

\begin{equation}
    g
    =
    \Psi\circ\mathcal G
\end{equation}

の

\begin{equation}
    \text{非負性},
    \qquad
    \text{可測性},
    \qquad
    \text{積分可能性},
    \qquad
    \text{正規化可能性}
\end{equation}

を扱う。

したがって、

\begin{equation}
    \boxed{
    \begin{array}{c}
        \text{Solver の集合論的・構造論的定義}
        \\[2mm]
        \Downarrow
        \\[2mm]
        \text{一般評価空間の構成}
        \\[2mm]
        \Downarrow
        \\[2mm]
        \text{可測構造の導入}
        \\[2mm]
        \Downarrow
        \\[2mm]
        \text{評価関数の測度論的性質}
    \end{array}
    }
\end{equation}

という順序で理解される。


%------------------------------------------------------------
\subsection{(C.16) 本節における証明事項と仮定事項の区別}
\label{subsec:C16}

本節では、以下を明確に区別する。

\begin{enumerate}
    \item
    $\mathcal G$ および $\Psi$ が可測であるという仮定から、
    \[
        g=\Psi\circ\mathcal G
    \]
    が可測であること。

    \item
    \[
        0<\int_{\Solver}g\,d\mu_{\Solver}<\infty
    \]
    が成立するならば、
    \[
        \widehat g=\frac{g}{Z}
    \]
    が正規化されること。

    \item
    ECA の具体的評価構造において、
    \[
        B\geq0,\qquad I_S\geq0,\qquad K\geq0
    \]
    が保証されるならば、
    \[
        g=BI_SK\geq0
    \]
    が成立すること。

    \item
    $\mathcal G$ の存在そのもの、
    $\Psi$ の可測性そのもの、
    \[
        0<
        \int_{\Solver}g\,d\mu_{\Solver}
        <
        \infty
    \]
    の成立そのものは、各構造から別途検証されるべき事項であり、
    定義だけから自動的に導かれるものではない。
\end{enumerate}

したがって、(C) は

\begin{equation}
    \boxed{
    \text{測度論的性質を定義すること}
    \quad\text{と}\quad
    \text{それらの成立を証明すること}
    }
\end{equation}

を混同しない構成とする。

特に、正規化可能性については

\begin{equation}
    0<
    \int_{\Solver}g\,d\mu_{\Solver}
    <
    \infty
\end{equation}

を成立条件として明示し、この条件が ECA の具体的な
Solver 集合に対して成立するかどうかは、後続の各定理において
個別に検証する。

%################################################################
% (D)
%################################################################

\section{(D) 一般解の存在}
\label{sec:D}

%------------------------------------------------------------
\subsection{(D.1) 一般解の存在問題}
\label{subsec:D1}

(B) において、Solver 集合 $\Solver$ に対する一般解を

\begin{equation}
    \mathcal{G}:
    \Solver
    \longrightarrow
    X
    \label{eq:D-general-solution-map}
\end{equation}

という型を持つ写像として定めた。

ここで重要なのは、写像の型を定めることと、
その型を満たす写像が実際に存在することは別の問題であるという点である。

したがって (D) では、

\begin{equation}
    \boxed{
        \exists\,
        \mathcal{G}:
        \Solver
        \longrightarrow
        X
    }
    \label{eq:D-existence}
\end{equation}

を一般解の存在命題として扱う。

ここでいう存在とは、各 $S\in\Solver$ に対して個別に
何らかの評価値を割り当てることではない。

すなわち、Solver 集合全体

\begin{equation}
    \Solver
\end{equation}

に属する任意の Solver に対して、その Solver の評価構造を
同一の規則によって $X$ の元へ対応させる写像

\begin{equation}
    \mathcal{G}
\end{equation}

が存在することを意味する。

したがって、

\begin{equation}
    \forall S\in\Solver,
    \qquad
    \mathcal{G}(S)\in X
    \label{eq:D-pointwise-existence}
\end{equation}

が成立する。


%------------------------------------------------------------
\subsection{(D.2) 一般解と評価構造}
\label{subsec:D2}

(B) において定義された一般評価空間を

\begin{equation}
    X
\end{equation}

とする。

一般解

\begin{equation}
    \mathcal{G}:
    \Solver\to X
\end{equation}

が存在するならば、任意の

\begin{equation}
    S\in\Solver
\end{equation}

について

\begin{equation}
    \mathcal{G}(S)
    \in X
\end{equation}

である。

ここで $\mathcal{G}(S)$ は、Solver $S$ に対して
(B) で定義された評価構造を適用した結果として得られる
一般評価を表す。

したがって、一般解は単なる数値を返す写像ではなく、

\begin{equation}
    \boxed{
        \mathcal{G}(S)
        =
        \text{Solver $S$ に対応する一般評価構造}
    }
\end{equation}

として解釈する。


%------------------------------------------------------------
\subsection{(D.3) 一般解から最終評価への写像}
\label{subsec:D3}

(B) において、評価空間から実数値を取り出す最終評価写像を

\begin{equation}
    \Psi:
    X
    \longrightarrow
    \R
    \label{eq:D-final-evaluation-map}
\end{equation}

とした。

一般解 $\mathcal{G}$ が存在するとき、

\begin{equation}
    g
    :=
    \Psi\circ\mathcal{G}
    \label{eq:D-final-evaluation}
\end{equation}

によって Solver 集合上の最終評価関数

\begin{equation}
    g:
    \Solver
    \longrightarrow
    \R
\end{equation}

を得ることができる。

すなわち、任意の $S\in\Solver$ に対して

\begin{equation}
    \boxed{
        g(S)
        =
        \Psi\!\left(\mathcal{G}(S)\right)
    }
    \label{eq:D-final-evaluation-explicit}
\end{equation}

である。

したがって、一般解の存在

\begin{equation}
    \mathcal{G}:\Solver\to X
\end{equation}

および最終評価写像

\begin{equation}
    \Psi:X\to\R
\end{equation}

が与えられれば、

\begin{equation}
    g=\Psi\circ\mathcal{G}
\end{equation}

が定義される。


%------------------------------------------------------------
\subsection{(D.4) 一般解の存在と (A)--(C) の関係}
\label{subsec:D4}

(A) では Solver 集合およびその集合論的・構造論的基礎を定めた。

(B) では、その Solver 集合を評価するための一般評価空間

\begin{equation}
    X
\end{equation}

および一般解の型

\begin{equation}
    \mathcal{G}:\Solver\to X
\end{equation}

を定めた。

(C) では、評価空間および Solver 集合に対して必要となる
可測構造、測度、積分可能性、正規化可能性等の
測度論的性質を定めた。

したがって、(D) における一般解は、

\begin{equation}
    \boxed{
    \Solver
    \xrightarrow{\mathcal{G}}
    X
    \xrightarrow{\Psi}
    \R
    }
    \label{eq:D-structure}
\end{equation}

という構造の中央に位置する。

ここで、

\begin{equation}
    \mathcal{G}
\end{equation}

が Solver の構造を一般評価空間 $X$ に対応させ、

\begin{equation}
    \Psi
\end{equation}

がその一般評価構造から最終的な評価値を取り出す。


%------------------------------------------------------------
\subsection{(D.5) 一般解の存在に関する命題}
\label{subsec:D5}

\begin{proposition}[一般解の存在]
\label{prop:D-general-solution-existence}

Solver 集合 $\Solver$ と一般評価空間 $X$ が (A)--(B) において
定義されているとする。

さらに、各

\begin{equation}
    S\in\Solver
\end{equation}

に対して、その評価構造を一意に定める規則

\begin{equation}
    S
    \longmapsto
    x_S
    \in X
\end{equation}

が与えられているならば、

\begin{equation}
    \mathcal{G}(S):=x_S
\end{equation}

によって

\begin{equation}
    \mathcal{G}:
    \Solver
    \longrightarrow
    X
\end{equation}

を定義できる。
\end{proposition}

\begin{proof}

仮定より、任意の

\begin{equation}
    S\in\Solver
\end{equation}

に対して、対応する評価構造

\begin{equation}
    x_S\in X
\end{equation}

が一意に定まる。

そこで、

\begin{equation}
    \mathcal{G}(S):=x_S
\end{equation}

と定義する。

この定義により、

\begin{equation}
    \forall S\in\Solver,
    \qquad
    \mathcal{G}(S)\in X
\end{equation}

が成立する。

したがって、

\begin{equation}
    \mathcal{G}:
    \Solver\to X
\end{equation}

は well-defined な写像である。

よって、

\begin{equation}
    \exists\,
    \mathcal{G}:
    \Solver\to X
\end{equation}

が成立する。
\end{proof}


%------------------------------------------------------------
\subsection{(D.6) 一般解の存在と最終評価関数}
\label{subsec:D6}

\begin{proposition}[一般解から最終評価関数の存在が従う]
\label{prop:D-final-evaluation-existence}

一般解

\begin{equation}
    \mathcal{G}:
    \Solver\to X
\end{equation}

および最終評価写像

\begin{equation}
    \Psi:
    X\to\R
\end{equation}

が存在するとする。

このとき、

\begin{equation}
    g
    :=
    \Psi\circ\mathcal{G}
\end{equation}

によって、

\begin{equation}
    g:
    \Solver\to\R
\end{equation}

が存在する。
\end{proposition}

\begin{proof}

$\mathcal{G}$ は

\begin{equation}
    \mathcal{G}:\Solver\to X
\end{equation}

であり、$\Psi$ は

\begin{equation}
    \Psi:X\to\R
\end{equation}

である。

したがって、写像の合成

\begin{equation}
    \Psi\circ\mathcal{G}
\end{equation}

が定義できる。

実際、任意の $S\in\Solver$ に対して

\begin{equation}
    \mathcal{G}(S)\in X
\end{equation}

であり、$\Psi$ の定義域は $X$ なので、

\begin{equation}
    \Psi(\mathcal{G}(S))
\end{equation}

が定義される。

よって、

\begin{equation}
    g(S)
    :=
    \Psi(\mathcal{G}(S))
\end{equation}

によって

\begin{equation}
    g:\Solver\to\R
\end{equation}

が定義される。
\end{proof}


%------------------------------------------------------------
\subsection{(D.7) 可測構造を考慮した一般解}
\label{subsec:D7}

(C) において Solver 集合および評価空間に可測構造を導入した。

したがって、

\begin{equation}
    (\Solver,\Sigma_{\Solver})
\end{equation}

および

\begin{equation}
    (X,\Sigma_X)
\end{equation}

を考える。

このとき、一般解を測度論的評価へ接続するためには、

\begin{equation}
    \mathcal{G}:
    (\Solver,\Sigma_{\Solver})
    \longrightarrow
    (X,\Sigma_X)
\end{equation}

が可測であることを要求する。

さらに、

\begin{equation}
    \Psi:
    (X,\Sigma_X)
    \longrightarrow
    (\R,\mathcal B(\R))
\end{equation}

が可測ならば、(C) の命題より

\begin{equation}
    g
    =
    \Psi\circ\mathcal{G}
\end{equation}

も可測となる。

したがって、

\begin{equation}
    \boxed{
    \begin{array}{ccccc}
        (\Solver,\Sigma_{\Solver})
        &\xrightarrow{\mathcal{G}}&
        (X,\Sigma_X)
        &\xrightarrow{\Psi}&
        (\R,\mathcal B(\R))
        \\[2mm]
        &&
        \Downarrow
        &&
        \\[-1mm]
        &&
        g=\Psi\circ\mathcal{G}
        &&
    \end{array}
    }
\end{equation}

という測度論的評価構造が得られる。


%------------------------------------------------------------
\subsection{(D.8) 一般解の存在と一般解の一意性の区別}
\label{subsec:D8}

(D) において示されるのは、基本的には

\begin{equation}
    \exists\,
    \mathcal{G}:
    \Solver\to X
\end{equation}

という一般解の存在である。

これに対して、

\begin{equation}
    \exists!\,
    \mathcal{G}:
    \Solver\to X
\end{equation}

すなわち一般解の一意性は、存在命題とは別の問題である。

特に、異なる Solver の表現が存在する場合には、
同一の構造を表現する Solver 間で一般解がどのように対応するかを
別途考える必要がある。

この問題については、(E) において Solver 同型

\begin{equation}
    T:
    \Solver_1
    \longrightarrow
    \Solver_2
\end{equation}

と一般解との関係を扱う。


%------------------------------------------------------------
\subsection{(D.9) 一般解の存在に関する注意}
\label{subsec:D9}

ここで注意すべきなのは、(A)--(C) において定義された構造から、

\begin{equation}
    \exists\,
    \mathcal{G}:
    \Solver\to X
\end{equation}

が無条件に従うとは限らないことである。

(B) において

\begin{equation}
    \mathcal{G}:
    \Solver\to X
\end{equation}

という型を定めたことは、

\begin{equation}
    \mathcal{G}
\end{equation}

という写像が実際に存在することとは異なる。

したがって、本節では、

\begin{equation}
    \boxed{
        \text{一般解の型の定義}
        \neq
        \text{一般解の存在証明}
    }
\end{equation}

を明確に区別する。

一般解の存在を証明するためには、各 Solver $S$ に対して
対応する評価構造

\begin{equation}
    x_S\in X
\end{equation}

が存在し、かつその対応

\begin{equation}
    S\mapsto x_S
\end{equation}

が well-defined であることを示す必要がある。


%------------------------------------------------------------
\subsection{(D.10) (D) の位置づけ}
\label{subsec:D10}

以上より、(A)--(D) の構造は

\begin{equation}
    \boxed{
    \begin{array}{c}
        \text{(A) Solver 集合の基礎構造}
        \\[2mm]
        \Downarrow
        \\[2mm]
        \text{(B) 一般評価空間と一般解の型}
        \\[2mm]
        \Downarrow
        \\[2mm]
        \text{(C) 測度論的評価構造}
        \\[2mm]
        \Downarrow
        \\[2mm]
        \text{(D) 一般解の存在}
    \end{array}
    }
\end{equation}

という順序で構成される。

(D) における中心命題は、

\begin{equation}
    \boxed{
        \exists\,
        \mathcal{G}:
        \Solver\to X
    }
\end{equation}

である。

この一般解が存在するとき、

\begin{equation}
    g
    =
    \Psi\circ\mathcal{G}
\end{equation}

によって Solver 集合全体に対する最終評価関数を構成できる。

一方、

\begin{equation}
    \exists!\,
    \mathcal{G}
\end{equation}

という一意性については (D) では仮定せず、
Solver 同型との関係を含めて (E) において扱う。

したがって、

\begin{equation}
    \boxed{
        \text{(D) 一般解の存在}
        \longrightarrow
        \text{(E) 一般解の構造的一意性}
    }
\end{equation}

という論理的順序を維持する。

%################################################################
% (E)
%################################################################

\section{(E) Solver 同型と一般解の構造的一意性}
\label{sec:E}

%------------------------------------------------------------
\subsection{(E.1) 本節の目的}
\label{subsec:E1}

(D) において、Solver 集合 $\Solver$ に対する一般解

\begin{equation}
    \mathcal{G}:
    \Solver
    \longrightarrow
    X
\end{equation}

の存在を扱った。

しかし、一般解が存在することだけでは、
その一般解が Solver の具体的な表現方法に依存しないことまでは
保証されない。

例えば、同一の Solver 構造を異なる集合

\begin{equation}
    \Solver_1,
    \qquad
    \Solver_2
\end{equation}

によって表現する場合、それぞれに対して

\begin{equation}
    \Phi_1:
    \Solver_1\to X_1,
    \qquad
    \Phi_2:
    \Solver_2\to X_2
\end{equation}

という評価写像を構成することが考えられる。

このとき重要となるのは、

\begin{equation}
    \text{「同一の Solver 構造に対して、
    一般解が同一の評価構造を与えるか」}
\end{equation}

という問題である。

(E) では、この性質を Solver 同型との可換性によって定式化する。


%------------------------------------------------------------
\subsection{(E.2) Solver 同型}
\label{subsec:E2}

二つの Solver 集合

\begin{equation}
    \Solver_1,
    \qquad
    \Solver_2
\end{equation}

を考える。

これらが Solver として同一の構造を持つことを、

\begin{equation}
    \boxed{
        T:
        \Solver_1
        \overset{\sim}{\longrightarrow}
        \Solver_2
    }
    \label{eq:E-solver-isomorphism}
\end{equation}

という Solver 同型によって表す。

ここで $T$ は単なる集合間の全単射としてではなく、
(A) において定められた Solver の構造を保存する写像として扱う。

したがって、

\begin{equation}
    T
    \text{ が Solver 同型}
\end{equation}

であることには、Solver の表現に依存しない
構造的対応が成立していることを含める。

以降では、$T$ によって対応する

\begin{equation}
    S\in\Solver_1
\end{equation}

と

\begin{equation}
    T(S)\in\Solver_2
\end{equation}

を同一の Solver 構造を異なる表現によって記述したものとして扱う。


%------------------------------------------------------------
\subsection{(E.3) 一般解の構造的一意性}
\label{subsec:E3}

それぞれの Solver 集合について一般解

\begin{equation}
    \Phi_1:
    \Solver_1\to X_1,
    \qquad
    \Phi_2:
    \Solver_2\to X_2
\end{equation}

を考える。

まず、評価空間が共通して

\begin{equation}
    X_1=X_2=X
\end{equation}

である場合を考える。

このとき、一般解が Solver の構造によって一意に定まるとは、

\begin{equation}
    \boxed{
        \Phi_2\circ T
        =
        \Phi_1
    }
    \label{eq:E-strong-uniqueness}
\end{equation}

が成立することと定める。

すなわち、任意の

\begin{equation}
    S\in\Solver_1
\end{equation}

に対して、

\begin{equation}
    \boxed{
        \Phi_2(T(S))
        =
        \Phi_1(S)
    }
    \label{eq:E-pointwise-correspondence}
\end{equation}

が成立する。

これは、Solver の具体的な表現を変更しても、
Solver 同型によって対応する Solver に対して
同一の一般評価構造が得られることを意味する。

したがって、ここでいう一意性は、

\begin{equation}
    \Phi_1=\Phi_2
\end{equation}

という意味ではなく、

\begin{equation}
    \boxed{
        \Phi_2\circ T=\Phi_1
    }
\end{equation}

という Solver 同型を介した一意性である。


%------------------------------------------------------------
\subsection{(E.4) 評価写像の可換性}
\label{subsec:E4}

一般解を構成する評価写像

\begin{equation}
    \Phi_1:
    \Solver_1\to X_1,
    \qquad
    \Phi_2:
    \Solver_2\to X_2
\end{equation}

および Solver 同型

\begin{equation}
    T:
    \Solver_1
    \overset{\sim}{\longrightarrow}
    \Solver_2
\end{equation}

を考える。

評価空間が共通する場合、

\begin{equation}
    X_1=X_2=X
\end{equation}

ならば、

\begin{equation}
    \boxed{
        \begin{CD}
            \Solver_1 @>{T}>> \Solver_2\\
            @V{\Phi_1}VV @VV{\Phi_2}V\\
            X @= X
        \end{CD}
    }
    \label{eq:E-commutative-diagram}
\end{equation}

という可換図式によって一般解の構造的一意性を表現できる。

すなわち、

\begin{equation}
    \Phi_2\circ T
    =
    \Phi_1
\end{equation}

である。

この関係は、Solver の表現と一般解の評価との間に
矛盾が生じないことを意味する。


%------------------------------------------------------------
\subsection{(E.5) 評価空間が異なる場合}
\label{subsec:E5}

一般には、

\begin{equation}
    X_1\neq X_2
\end{equation}

であることも許容する。

この場合、単に

\begin{equation}
    \Phi_2\circ T=\Phi_1
\end{equation}

と書くことはできない。

そこで、評価空間間の対応写像

\begin{equation}
    U:
    X_1
    \longrightarrow
    X_2
\end{equation}

を導入する。

このとき、評価構造の対応を

\begin{equation}
    \boxed{
        U\circ\Phi_1
        =
        \Phi_2\circ T
    }
    \label{eq:E-evaluation-space-commutativity}
\end{equation}

によって定める。

すなわち、任意の

\begin{equation}
    S\in\Solver_1
\end{equation}

に対して、

\begin{equation}
    U(\Phi_1(S))
    =
    \Phi_2(T(S))
\end{equation}

である。

この関係により、評価空間そのものが異なる場合でも、
Solver 同型によって対応する評価構造を比較することができる。


%------------------------------------------------------------
\subsection{(E.6) 評価構造を保存する対応}
\label{subsec:E6}

ここで、$U:X_1\to X_2$ は単なる集合間の写像ではなく、
ECA において定義された評価構造を保存する対応として扱う。

すなわち、

\begin{equation}
    U:
    X_1
    \longrightarrow
    X_2
\end{equation}

が評価構造を保存するとは、評価空間 $X_1$ 上で定義された
構造と、$U$ によって $X_2$ に移された構造とが
ECA の評価規則に関して対応することをいう。

したがって、

\begin{equation}
    \boxed{
        \text{評価空間の対応}
        \neq
        \text{単なる集合の同型}
    }
\end{equation}

である。

必要に応じて $U$ は、評価空間に導入された構造
および (C) において定義された測度論的構造を保存する対応として
扱う。

ただし、具体的にどの構造を保存する必要があるかについては、
各評価空間の構成に応じて個別に定める。


%------------------------------------------------------------
\subsection{(E.7) Representation Independence}
\label{subsec:E7}

\begin{definition}[Representation Independence]
\label{def:E-representation-independence}

Solver 集合として同一の構造を表す二つの表現

\begin{equation}
    \Solver_1,
    \qquad
    \Solver_2
\end{equation}

に対して Solver 同型

\begin{equation}
    T:
    \Solver_1
    \overset{\sim}{\longrightarrow}
    \Solver_2
\end{equation}

が存在し、それぞれの一般解

\begin{equation}
    \Phi_1:\Solver_1\to X_1,
    \qquad
    \Phi_2:\Solver_2\to X_2
\end{equation}

について、評価空間間の構造保存写像

\begin{equation}
    U:X_1\to X_2
\end{equation}

が存在して

\begin{equation}
    U\circ\Phi_1
    =
    \Phi_2\circ T
\end{equation}

を満たすとき、これらの評価構造は
\emph{Representation Independent} であるという。
\end{definition}

特に、

\begin{equation}
    X_1=X_2=X
\end{equation}

かつ

\begin{equation}
    U=\id_X
\end{equation}

の場合には、

\begin{equation}
    \Phi_2\circ T
    =
    \Phi_1
\end{equation}

に帰着する。


%------------------------------------------------------------
\subsection{(E.8) 最終評価関数への帰結}
\label{subsec:E8}

(C) において、一般解から最終評価関数を

\begin{equation}
    g_i
    :=
    \Psi_i\circ\Phi_i,
    \qquad
    i=1,2
\end{equation}

と定義する。

すなわち、

\begin{equation}
    g_1
    =
    \Psi_1\circ\Phi_1,
    \qquad
    g_2
    =
    \Psi_2\circ\Phi_2
\end{equation}

である。

ここで、評価空間間の対応

\begin{equation}
    U:X_1\to X_2
\end{equation}

が最終評価を保存することを、

\begin{equation}
    \boxed{
        \Psi_2\circ U
        =
        \Psi_1
    }
    \label{eq:E-final-evaluation-preservation}
\end{equation}

と定める。

さらに、

\begin{equation}
    U\circ\Phi_1
    =
    \Phi_2\circ T
\end{equation}

が成立するとする。

このとき、

\begin{align}
    g_2\circ T
    &=
    \Psi_2\circ\Phi_2\circ T
    \\
    &=
    \Psi_2\circ U\circ\Phi_1
    \\
    &=
    \Psi_1\circ\Phi_1
    \\
    &=
    g_1.
\end{align}

したがって、

\begin{equation}
    \boxed{
        g_2\circ T
        =
        g_1
    }
    \label{eq:E-final-evaluation-invariance}
\end{equation}

が成立する。


%------------------------------------------------------------
\subsection{(E.9) 一般解の構造的一意性定理}
\label{subsec:E9}

\begin{theorem}[Solver 構造に基づく一般解の構造的一意性]
\label{thm:E-general-solution-uniqueness}

Solver 集合

\begin{equation}
    \Solver_1,
    \qquad
    \Solver_2
\end{equation}

および Solver 同型

\begin{equation}
    T:
    \Solver_1
    \overset{\sim}{\longrightarrow}
    \Solver_2
\end{equation}

を考える。

各 Solver 集合に対して一般解

\begin{equation}
    \Phi_1:
    \Solver_1\to X_1,
    \qquad
    \Phi_2:
    \Solver_2\to X_2
\end{equation}

が与えられているとする。

さらに、評価空間間の構造保存写像

\begin{equation}
    U:
    X_1\to X_2
\end{equation}

が存在し、

\begin{equation}
    U\circ\Phi_1
    =
    \Phi_2\circ T
    \label{eq:E-theorem-hypothesis-1}
\end{equation}

を満たすとする。

このとき、二つの一般解は Solver 同型 $T$ および
評価空間の対応 $U$ に関して同一の評価構造を表す。

さらに、最終評価写像

\begin{equation}
    \Psi_1:X_1\to\R,
    \qquad
    \Psi_2:X_2\to\R
\end{equation}

が

\begin{equation}
    \Psi_2\circ U
    =
    \Psi_1
    \label{eq:E-theorem-hypothesis-2}
\end{equation}

を満たすならば、

\begin{equation}
    g_2\circ T
    =
    g_1
\end{equation}

が成立する。
\end{theorem}


%------------------------------------------------------------
\subsection{(E.10) 定理の証明}
\label{subsec:E10}

\begin{proof}

まず、

\begin{equation}
    U\circ\Phi_1
    =
    \Phi_2\circ T
\end{equation}

が仮定されている。

したがって、任意の

\begin{equation}
    S\in\Solver_1
\end{equation}

に対して、

\begin{equation}
    U(\Phi_1(S))
    =
    \Phi_2(T(S))
\end{equation}

が成立する。

よって、$T$ によって対応する Solver は、
$U$ によって対応する評価構造を与える。

これは一般解が Solver の具体的表現そのものではなく、
Solver の構造に従って定まることを意味する。

次に最終評価関数について考える。

定義より、

\begin{equation}
    g_1
    =
    \Psi_1\circ\Phi_1,
    \qquad
    g_2
    =
    \Psi_2\circ\Phi_2
\end{equation}

である。

したがって、

\begin{align}
    g_2\circ T
    &=
    \Psi_2\circ\Phi_2\circ T
    \\
    &=
    \Psi_2\circ U\circ\Phi_1
    \\
    &=
    \Psi_1\circ\Phi_1
    \\
    &=
    g_1.
\end{align}

ここで第2の等号には

\begin{equation}
    \Phi_2\circ T
    =
    U\circ\Phi_1
\end{equation}

を用い、第3の等号には

\begin{equation}
    \Psi_2\circ U
    =
    \Psi_1
\end{equation}

を用いた。

したがって、

\begin{equation}
    \boxed{
        g_2\circ T=g_1
    }
\end{equation}

が成立する。
\end{proof}


%------------------------------------------------------------
\subsection{(E.11) 「一意性」の意味}
\label{subsec:E11}

ここでいう「一般解の構造的一意性」は、異なる Solver 集合や
異なる評価空間が文字通り同一であることを意味しない。

すなわち、

\begin{equation}
    \Solver_1=\Solver_2
\end{equation}

や

\begin{equation}
    X_1=X_2
\end{equation}

を要求するものではない。

むしろ、

\begin{equation}
    T:
    \Solver_1
    \overset{\sim}{\longrightarrow}
    \Solver_2
\end{equation}

および

\begin{equation}
    U:
    X_1\to X_2
\end{equation}

によって対応関係が与えられ、

\begin{equation}
    U\circ\Phi_1
    =
    \Phi_2\circ T
\end{equation}

が成立することによって、

\begin{equation}
    \boxed{
        \text{異なる表現から得られた一般解が
        同一の評価構造を表す}
    }
\end{equation}

ことを意味する。

したがって、本節における一意性は、

\begin{equation}
    \boxed{
        \text{Solver 構造に対する
        一意性}
    }
\end{equation}

であり、単なる写像の値の一致ではない。


%------------------------------------------------------------
\subsection{(E.12) 評価空間が共通の場合}
\label{subsec:E12}

特に、

\begin{equation}
    X_1=X_2=X
\end{equation}

かつ

\begin{equation}
    U=\id_X
\end{equation}

である場合には、定理の条件は

\begin{equation}
    \Phi_2\circ T
    =
    \Phi_1
\end{equation}

へ簡約される。

したがって、

\begin{equation}
    \boxed{
        \Phi_2(T(S))
        =
        \Phi_1(S)
        \qquad
        (\forall S\in\Solver_1)
    }
\end{equation}

である。

さらに、

\begin{equation}
    \Psi_1=\Psi_2=\Psi
\end{equation}

ならば、

\begin{align}
    g_2\circ T
    &=
    \Psi\circ\Phi_2\circ T
    \\
    &=
    \Psi\circ\Phi_1
    \\
    &=
    g_1.
\end{align}

したがって、

\begin{equation}
    \boxed{
        g_2\circ T=g_1
    }
\end{equation}

となる。


%------------------------------------------------------------
\subsection{(E.13) (B)--(D) との関係}
\label{subsec:E13}

(B) において一般解を

\begin{equation}
    \mathcal{G}:
    \Solver\to X
\end{equation}

という型を持つ写像として定義した。

(C) では、その一般解から得られる最終評価

\begin{equation}
    g
    =
    \Psi\circ\mathcal{G}
\end{equation}

について測度論的性質を定めた。

(D) では、

\begin{equation}
    \exists\,
    \mathcal{G}:
    \Solver\to X
\end{equation}

という一般解の存在を扱った。

これに対して (E) では、異なる Solver 表現に対して得られる
一般解が Solver 同型を通じて対応することを扱う。

したがって、

\begin{equation}
    \boxed{
    \begin{array}{c}
        \text{(B) 一般解の型}\\
        \Downarrow\\
        \text{(C) 最終評価の測度論的性質}\\
        \Downarrow\\
        \text{(D) 一般解の存在}\\
        \Downarrow\\
        \text{(E) 一般解の構造的一意性}
    \end{array}
    }
\end{equation}

という論理的順序になる。


%------------------------------------------------------------
\subsection{(E.14) 後続する (F)--(H) への接続}
\label{subsec:E14}

(E) において得られた構造的一意性は、
Solver の表現方法が異なる場合にも一般解の評価構造を
対応させるための基礎となる。

特に、離散 Solver と連続 Solver がそれぞれ

\begin{equation}
    \Phi_{\mathrm{disc}}
    :
    \Solver_{\mathrm{disc}}
    \to
    X_{\mathrm{disc}}
\end{equation}

および

\begin{equation}
    \Phi_{\mathrm{cont}}
    :
    \Solver_{\mathrm{cont}}
    \to
    X_{\mathrm{cont}}
\end{equation}

として構成される場合、それらが Solver 同型および
評価空間の構造保存写像によって対応するならば、

\begin{equation}
    U\circ\Phi_{\mathrm{disc}}
    =
    \Phi_{\mathrm{cont}}\circ T
\end{equation}

という関係によって両者の一般解構造を比較できる。

この構造を基礎として、

\begin{equation}
    \text{(F) 離散評価}
    \longrightarrow
    \text{(G) 連続評価}
    \longrightarrow
    \text{(H) 中間 Solver}
\end{equation}

において、異なる Solver 表現を統一的に扱う。


%------------------------------------------------------------
\subsection{(E.15) 本節における証明事項と仮定事項の区別}
\label{subsec:E15}

本節では、以下の事項を明確に区別する。

\begin{enumerate}

    \item
    Solver 同型
    \[
        T:\Solver_1\overset{\sim}{\longrightarrow}\Solver_2
    \]
    が与えられた場合に、

    \[
        U\circ\Phi_1
        =
        \Phi_2\circ T
    \]

    が成立することは、一般解が Solver の構造に対して
    一貫していることを表す条件である。

    \item
    上記の可換性から、

    \[
        \Psi_2\circ U=\Psi_1
    \]

    を仮定すれば、

    \[
        g_2\circ T=g_1
    \]

    が合成写像の計算によって導かれる。

    \item
    一方、

    \[
        U\circ\Phi_1
        =
        \Phi_2\circ T
    \]

    そのものが任意の Solver に対して
    (A)--(D) だけから自動的に成立するとは限らない。

    \item
    同様に、

    \[
        \Psi_2\circ U=\Psi_1
    \]

    も最終評価写像および評価空間の構造に応じて
    別途確認されるべき条件である。

\end{enumerate}

したがって、(E) の中心的な数学的帰結は、

\begin{equation}
    \boxed{
        \left.
        \begin{array}{c}
            U\circ\Phi_1
            =
            \Phi_2\circ T
            \\[2mm]
            \Psi_2\circ U
            =
            \Psi_1
        \end{array}
        \right\}
        \Longrightarrow
        g_2\circ T=g_1
    }
\end{equation}

である。


%------------------------------------------------------------
\subsection{(E.16) 本節の結論}
\label{subsec:E16}

以上により、(E) における一般解の構造的一意性は、

\begin{equation}
    \boxed{
        U\circ\Phi_1
        =
        \Phi_2\circ T
    }
\end{equation}

という評価構造の可換性によって定式化される。

さらに最終評価を保存する条件

\begin{equation}
    \boxed{
        \Psi_2\circ U
        =
        \Psi_1
    }
\end{equation}

が成立すれば、

\begin{equation}
    \boxed{
        g_2\circ T
        =
        g_1
    }
\end{equation}

が従う。

したがって、

\begin{equation}
    \boxed{
    \begin{array}{c}
        \text{Solver 同型}\\
        T:\Solver_1\overset{\sim}{\longrightarrow}\Solver_2
        \\[2mm]
        \Downarrow
        \\[2mm]
        \text{評価構造の可換性}\\
        U\circ\Phi_1=\Phi_2\circ T
        \\[2mm]
        \Downarrow
        \\[2mm]
        \text{一般解の構造的一意性}
        \\[2mm]
        \Downarrow
        \\[2mm]
        \text{最終評価の表現独立性}\\
        g_2\circ T=g_1
    \end{array}
    }
\end{equation}

という構造が得られる。

ここで重要なのは、一般解の「一意性」を
異なる Solver や評価空間そのものを同一視することによって
定義していない点である。

むしろ、

\begin{equation}
    \boxed{
        \text{Solver 構造を保存する対応によって
        評価構造が可換する}
    }
\end{equation}

ことをもって一般解の構造的一意性とする。

この定式化により、(F)--(H) において異なる Solver 表現を
同一の一般解構造のもとで扱うための基礎が与えられる。

%################################################################
% (F)
%################################################################

\section{(F) 離散 Solver の細分化極限}
\label{sec:F-discrete-limit}

%------------------------------------------------------------
\subsection{(F.1) 本節の目的}
\label{subsec:F1}

本節では、(A)--(E) において構成された Solver 集合、
一般評価空間、評価写像、可測構造および正規化測度を用いて、
離散 Solver による一般解と連続形式の一般解との関係を扱う。

特に、有限段階 $N$ において構成される離散評価

\begin{equation}
    T_N
\end{equation}

を考え、その細分化極限

\begin{equation}
    N\to\infty
\end{equation}

において対応する測度積分へ収束するための条件を定式化する。

本節では新たな正規化を導入しない。

(A)--(D) において既に定義された重み関数、
可測構造および正規化測度を用いて、離散 Solver の評価を記述する。

ECA において既に与えられている離散形式は

\begin{equation}
    T_{\mathrm{overlay},N}
    =
    \frac{1}{N}
    \sum_{j=1}^{N}
    B_j I_j K_j
    \label{eq:F-original-discrete}
\end{equation}

であり、対応する連続形式は

\begin{equation}
    T_{\mathrm{overlay}}
    =
    \int
    B(x)I(x)K(x)\,d\mu(x)
    \label{eq:F-original-continuous}
\end{equation}

で表される。

したがって、本節では既存の離散式を、
(A)--(D) において定義された測度構造に基づく
有限段階の評価として一般化する。


%------------------------------------------------------------
\subsection{(F.2) Solver 集合と評価構造}
\label{subsec:F2}

(A)--(D) において定義された Solver 集合を

\begin{equation}
    \mathcal S
    =
    \{S\mid S=(\mathcal A_S,\mathcal R_S)\}
    \label{eq:F-solver-set}
\end{equation}

とする。

各

\begin{equation}
    S\in\mathcal S
\end{equation}

は、公理選択および推論規則・構成制約によって定まる Solver
として扱う。

(B)--(D) において構成された一般評価空間を

\begin{equation}
    X
\end{equation}

とし、Solver 集合から評価空間への一般評価写像を

\begin{equation}
    \Phi:
    \mathcal S
    \longrightarrow
    X
    \label{eq:F-general-evaluation-map}
\end{equation}

とする。

また、ECA の評価に現れる三つの量を

\begin{equation}
    B,\qquad I,\qquad K
\end{equation}

とし、その積を

\begin{equation}
    f(x)
    :=
    B(x)I(x)K(x)
    \label{eq:F-integrand}
\end{equation}

と定義する。

ここで $f$ は、離散 Solver の各評価要素を
連続的な評価空間上の関数として表現するために用いる。


%------------------------------------------------------------
\subsection{(F.3) 正規化測度と細分化}
\label{subsec:F3}

(A)--(D) において定義された可測空間を

\begin{equation}
    (I,\mathcal F)
\end{equation}

とし、その正規化測度を

\begin{equation}
    \mu:\mathcal F\to[0,1]
\end{equation}

とする。

正規化条件として

\begin{equation}
    \mu(I)=1
    \label{eq:F-measure-normalization}
\end{equation}

を満たすものとする。

有限段階 $N$ における Solver 集合の細分化を

\begin{equation}
    \mathcal P_N
    =
    \left\{
        E_1^{(N)},
        E_2^{(N)},
        \ldots,
        E_N^{(N)}
    \right\}
    \label{eq:F-partition}
\end{equation}

とする。

各

\begin{equation}
    E_j^{(N)}
\end{equation}

は $\mathcal F$ の元であり、

\begin{equation}
    E_j^{(N)}\in\mathcal F
\end{equation}

を満たす。

さらに、

\begin{equation}
    E_p^{(N)}
    \cap
    E_q^{(N)}
    =
    \varnothing
    \qquad
    (p\neq q)
    \label{eq:F-disjoint}
\end{equation}

および

\begin{equation}
    \bigcup_{j=1}^{N}
    E_j^{(N)}
    =
    I
    \label{eq:F-cover}
\end{equation}

を満たすものとする。

各分割要素に対応する測度要素を

\begin{equation}
    \boxed{
        \Delta\mu_j^{(N)}
        :=
        \mu\!\left(E_j^{(N)}\right)
    }
    \label{eq:F-measure-element}
\end{equation}

と定義する。

測度の加法性と \eqref{eq:F-measure-normalization} より、

\begin{equation}
    \sum_{j=1}^{N}
    \Delta\mu_j^{(N)}
    =
    \sum_{j=1}^{N}
    \mu\!\left(E_j^{(N)}\right)
    =
    \mu(I)
    =
    1
    \label{eq:F-partition-normalization}
\end{equation}

が成立する。


%------------------------------------------------------------
\subsection{(F.4) 離散 Solver の評価}
\label{subsec:F4}

各分割要素

\begin{equation}
    E_j^{(N)}
\end{equation}

に対して、その評価空間上の代表元を

\begin{equation}
    x_j^{(N)}\in X
    \label{eq:F-representative}
\end{equation}

とする。

このとき、各離散 Solver 要素における評価を

\begin{equation}
    B_j I_j K_j
    =
    f\!\left(x_j^{(N)}\right)
    \label{eq:F-discrete-evaluation}
\end{equation}

として表す。

ここで、

\begin{equation}
    f(x)=B(x)I(x)K(x)
\end{equation}

である。

したがって、有限段階 $N$ における離散一般解を

\begin{equation}
    \boxed{
        T_N
        =
        \sum_{j=1}^{N}
        f\!\left(x_j^{(N)}\right)
        \Delta\mu_j^{(N)}
    }
    \label{eq:F-discrete-general}
\end{equation}

と定義する。

\eqref{eq:F-discrete-evaluation} を代入すれば、

\begin{equation}
    T_N
    =
    \sum_{j=1}^{N}
    B_j I_j K_j
    \Delta\mu_j^{(N)}
    \label{eq:F-discrete-BIK}
\end{equation}

となる。


%------------------------------------------------------------
\subsection{(F.5) 一様正規化の場合}
\label{subsec:F5}

特に、一様正規化された有限段階の Solver 集合については、

\begin{equation}
    \boxed{
        \Delta\mu_j^{(N)}
        =
        \frac{1}{N}
    }
    \label{eq:F-uniform-measure}
\end{equation}

とする。

このとき、

\begin{equation}
    \sum_{j=1}^{N}
    \Delta\mu_j^{(N)}
    =
    \sum_{j=1}^{N}
    \frac1N
    =
    1
\end{equation}

であり、正規化条件を満たす。

したがって \eqref{eq:F-discrete-BIK} は

\begin{equation}
    \boxed{
        T_N
        =
        \frac1N
        \sum_{j=1}^{N}
        B_jI_jK_j
    }
    \label{eq:F-ECA-discrete}
\end{equation}

となる。

これは ECA において既に与えられている離散形式をそのまま回収する。


%------------------------------------------------------------
\subsection{(F.6) 細分化メッシュ}
\label{subsec:F6}

離散 Solver 集合の細分化の度合いを、
各分割要素の測度によって

\begin{equation}
    \boxed{
        \left\|
        \mathcal P_N
        \right\|_\mu
        :=
        \max_{1\leq j\leq N}
        \Delta\mu_j^{(N)}
    }
    \label{eq:F-measure-mesh}
\end{equation}

と定義する。

この量を測度メッシュと呼ぶ。

細分化極限を

\begin{equation}
    \boxed{
        \left\|
        \mathcal P_N
        \right\|_\mu
        \longrightarrow
        0
        \qquad
        (N\to\infty)
    }
    \label{eq:F-mesh-limit}
\end{equation}

と定義する。

これは単に

\begin{equation}
    N\to\infty
\end{equation}

となることだけを意味しない。

各分割要素に割り当てられた測度が

\begin{equation}
    \Delta\mu_j^{(N)}
    =
    \mu\!\left(E_j^{(N)}\right)
\end{equation}

であることから、

\begin{equation}
    \left\|
    \mathcal P_N
    \right\|_\mu
    =
    \max_j
    \mu\!\left(E_j^{(N)}\right)
    \longrightarrow0
\end{equation}

となることを要求する。


%------------------------------------------------------------
\subsection{(F.7) 細分化列の存在}
\label{subsec:F7}

\begin{proposition}[細分化列の存在]
\label{prop:F-partition-existence}

(A)--(D) によって構成された可測空間 $(I,\mathcal F)$ に対して、
各 $N\in\mathbb N$ について

\begin{equation}
    \mathcal P_N
    =
    \left\{
        E_1^{(N)},
        \ldots,
        E_N^{(N)}
    \right\}
\end{equation}

という可測分割を考えることができるとする。

すなわち、

\begin{equation}
    E_j^{(N)}\in\mathcal F,
\end{equation}

\begin{equation}
    E_p^{(N)}
    \cap
    E_q^{(N)}
    =
    \varnothing
    \qquad(p\neq q),
\end{equation}

および

\begin{equation}
    \bigcup_{j=1}^{N}
    E_j^{(N)}
    =
    I
\end{equation}

を満たす。

このとき、これらの分割によって

\begin{equation}
    \Delta\mu_j^{(N)}
    =
    \mu(E_j^{(N)})
\end{equation}

を定義できる。
\end{proposition}

本命題は、細分化された Solver 集合を
可測分割として扱うための基礎を与える。


%------------------------------------------------------------
\subsection{(F.8) 離散一般解の極限の存在}
\label{subsec:F8}

\begin{proposition}[離散一般解の極限の存在]
\label{prop:F-limit-existence}

前命題による細分化列

\begin{equation}
    \{\mathcal P_N\}_{N=1}^{\infty}
\end{equation}

を考える。

さらに、

\begin{equation}
    \left\|
    \mathcal P_N
    \right\|_\mu
    \longrightarrow0
\end{equation}

が成立し、

\begin{equation}
    T_N
    =
    \sum_{j=1}^{N}
    f\!\left(x_j^{(N)}\right)
    \Delta\mu_j^{(N)}
\end{equation}

が対応する測度積分

\begin{equation}
    \int_I f(x)\,d\mu(x)
\end{equation}

へ収束するならば、

\begin{equation}
    \lim_{N\to\infty}T_N
\end{equation}

が存在する。
\end{proposition}

したがって、

\begin{equation}
    \boxed{
        \text{細分化条件}
        +
        \text{積分近似の収束条件}
        \Longrightarrow
        \text{離散一般解の極限の存在}
    }
\end{equation}

である。


%------------------------------------------------------------
\subsection{(F.9) 離散極限定理}
\label{subsec:F9}

\begin{theorem}[離散 Solver の細分化極限定理]
\label{thm:F-discrete-limit}

(A)--(D) によって構成された Solver 集合、評価空間、
可測構造および正規化測度を用いる。

有限段階 $N$ における離散一般解を

\begin{equation}
    T_N
    =
    \sum_{j=1}^{N}
    f\!\left(x_j^{(N)}\right)
    \Delta\mu_j^{(N)}
    \label{eq:F-theorem-discrete}
\end{equation}

とする。

ここで

\begin{equation}
    f(x)
    =
    B(x)I(x)K(x)
    \label{eq:F-theorem-integrand}
\end{equation}

であり、

\begin{equation}
    \sum_{j=1}^{N}
    \Delta\mu_j^{(N)}
    =
    1
\end{equation}

を満たすものとする。

さらに、

\begin{equation}
    \left\|
    \mathcal P_N
    \right\|_\mu
    \longrightarrow0
\end{equation}

が成立し、$f$ が対応する測度空間上で可積分であり、
当該細分化による積分近似が

\begin{equation}
    \lim_{N\to\infty}
    \sum_{j=1}^{N}
    f\!\left(x_j^{(N)}\right)
    \Delta\mu_j^{(N)}
    =
    \int_I f(x)\,d\mu(x)
    \label{eq:F-integral-convergence}
\end{equation}

を満たすならば、

\begin{equation}
    \boxed{
        \lim_{N\to\infty}T_N
        =
        \int_I f(x)\,d\mu(x)
    }
    \label{eq:F-limit}
\end{equation}

が成立する。

したがって、

\begin{equation}
    \boxed{
        \lim_{N\to\infty}
        \sum_{j=1}^{N}
        B\!\left(x_j^{(N)}\right)
        I\!\left(x_j^{(N)}\right)
        K\!\left(x_j^{(N)}\right)
        \Delta\mu_j^{(N)}
        =
        \int_I
        B(x)I(x)K(x)\,d\mu(x)
    }
    \label{eq:F-ECA-limit}
\end{equation}

である。
\end{theorem}


%------------------------------------------------------------
\subsection{(F.10) 離散極限定理の証明}
\label{subsec:F10}

\begin{proof}

定義より、

\begin{equation}
    \Delta\mu_j^{(N)}
    =
    \mu\!\left(E_j^{(N)}\right)
\end{equation}

である。

また、

\begin{equation}
    \mathcal P_N
    =
    \left\{
        E_1^{(N)},
        \ldots,
        E_N^{(N)}
    \right\}
\end{equation}

は $I$ の可測分割であるから、

\begin{equation}
    \sum_{j=1}^{N}
    \Delta\mu_j^{(N)}
    =
    \sum_{j=1}^{N}
    \mu\!\left(E_j^{(N)}\right)
    =
    \mu(I)
    =
    1
\end{equation}

が成立する。

各分割要素に対応する評価点を

\begin{equation}
    x_j^{(N)}
\end{equation}

とし、

\begin{equation}
    f(x)=B(x)I(x)K(x)
\end{equation}

と定義すると、

\begin{equation}
    T_N
    =
    \sum_{j=1}^{N}
    f\!\left(x_j^{(N)}\right)
    \Delta\mu_j^{(N)}
\end{equation}

となる。

ここで仮定より、

\begin{equation}
    \left\|
    \mathcal P_N
    \right\|_\mu
    =
    \max_j
    \Delta\mu_j^{(N)}
    \longrightarrow0.
\end{equation}

さらに、積分近似の収束条件

\begin{equation}
    \lim_{N\to\infty}
    \sum_{j=1}^{N}
    f\!\left(x_j^{(N)}\right)
    \Delta\mu_j^{(N)}
    =
    \int_I f(x)\,d\mu(x)
\end{equation}

を仮定している。

したがって、

\begin{equation}
    \lim_{N\to\infty}T_N
    =
    \int_I f(x)\,d\mu(x)
\end{equation}

が成立する。

最後に

\begin{equation}
    f(x)=B(x)I(x)K(x)
\end{equation}

を代入することで、

\begin{equation}
    \boxed{
        \lim_{N\to\infty}T_N
        =
        \int_I
        B(x)I(x)K(x)\,d\mu(x)
    }
\end{equation}

を得る。
\end{proof}


%------------------------------------------------------------
\subsection{(F.11) ECA の Solver 集合に対する細分化極限についての論考}
\label{subsec:F11}

前節までの定理および命題は、

\begin{equation}
    \left\|
    \mathcal P_N
    \right\|_\mu
    \longrightarrow0
\end{equation}

という条件を満たす細分化列について記述したものである。

ここで、

\begin{equation}
    \boxed{
        \text{細分化極限の条件を仮定すること}
    }
\end{equation}

と

\begin{equation}
    \boxed{
        \text{ECA の Solver 集合が実際に
        その条件を満たすこと}
    }
\end{equation}

とは区別しなければならない。

後者については、本節では新たな証明を与えず、
ECA において既に導入されている離散式および連続式との関係から
論考する。


%------------------------------------------------------------
\subsubsection{(F.11.1) 一様正規化の場合}
\label{subsubsec:F11-1}

最も単純な場合として、

\begin{equation}
    \Delta\mu_j^{(N)}
    =
    \frac1N
\end{equation}

とする。

このとき、

\begin{equation}
    \left\|
    \mathcal P_N
    \right\|_\mu
    =
    \max_j\frac1N
    =
    \frac1N
\end{equation}

である。

したがって、

\begin{equation}
    \lim_{N\to\infty}
    \left\|
    \mathcal P_N
    \right\|_\mu
    =
    \lim_{N\to\infty}
    \frac1N
    =
    0
\end{equation}

となる。

このとき、

\begin{equation}
    T_N
    =
    \frac1N
    \sum_{j=1}^{N}
    B_jI_jK_j
\end{equation}

となり、これは ECA の既存の離散式そのものである。

したがって、一様正規化された場合には、
離散 Solver の段階数を増加させることと、
測度メッシュを細分化することが一致する。


%------------------------------------------------------------
\subsubsection{(F.11.2) 一般の正規化測度の場合}
\label{subsubsec:F11-2}

一様正規化を仮定しない場合には、

\begin{equation}
    \Delta\mu_j^{(N)}
    =
    \mu\!\left(E_j^{(N)}\right)
\end{equation}

である。

したがって、

\begin{equation}
    \left\|
    \mathcal P_N
    \right\|_\mu
    =
    \max_j
    \mu\!\left(E_j^{(N)}\right)
\end{equation}

となる。

この場合、細分化極限の成立条件は

\begin{equation}
    \boxed{
        \lim_{N\to\infty}
        \max_j
        \mu\!\left(E_j^{(N)}\right)
        =
        0
    }
    \label{eq:F-general-mesh-condition}
\end{equation}

となる。

したがって、

\begin{equation}
    N\to\infty
\end{equation}

だけでは一般には十分ではない。

必要なのは、

\begin{equation}
    \max_j
    \mu\!\left(E_j^{(N)}\right)
    \longrightarrow0
\end{equation}

という測度的な細分化条件である。


%------------------------------------------------------------
\subsubsection{(F.11.3) 離散式から連続式への接続}
\label{subsubsec:F11-3}

以上をまとめると、離散形式

\begin{equation}
    T_N
    =
    \sum_{j=1}^{N}
    f\!\left(x_j^{(N)}\right)
    \Delta\mu_j^{(N)}
\end{equation}

に対して、

\begin{equation}
    \Delta\mu_j^{(N)}
    =
    \mu\!\left(E_j^{(N)}\right)
\end{equation}

および

\begin{equation}
    \left\|
    \mathcal P_N
    \right\|_\mu
    =
    \max_j
    \Delta\mu_j^{(N)}
    \longrightarrow0
\end{equation}

を考えることで、

\begin{equation}
    T_N
    \longrightarrow
    \int_I f(x)\,d\mu(x)
\end{equation}

という接続を得る。

さらに、

\begin{equation}
    f(x)=B(x)I(x)K(x)
\end{equation}

であるから、

\begin{equation}
    \boxed{
        T_N
        \longrightarrow
        \int_I
        B(x)I(x)K(x)\,d\mu(x)
    }
\end{equation}

となる。

したがって、ECA における離散 Solver から連続形式への移行は、

\begin{equation}
    \boxed{
    \text{離散 Solver}
    \longrightarrow
    \text{測度による細分化}
    \longrightarrow
    \text{メッシュの消失}
    \longrightarrow
    \text{連続評価}
    }
\end{equation}

という構造として理解できる。


%------------------------------------------------------------
\subsection{(F.12) (E) との関係}
\label{subsec:F12}

(E) において、Solver 同型に対して一般解の評価構造が
可換することを定式化した。

すなわち、

\begin{equation}
    U\circ\Phi_1
    =
    \Phi_2\circ T
\end{equation}

によって、異なる Solver 表現の間で
一般評価構造が対応する。

(F) では、その評価構造を有限段階において細分化し、

\begin{equation}
    \mathcal P_N
    =
    \left\{
        E_1^{(N)},
        \ldots,
        E_N^{(N)}
    \right\}
\end{equation}

として表現する。

したがって、(F) における離散一般解

\begin{equation}
    T_N
    =
    \sum_{j=1}^{N}
    f\!\left(x_j^{(N)}\right)
    \Delta\mu_j^{(N)}
\end{equation}

は、(E) において定義された一般評価構造を
有限段階の測度的分割によって表現したものと解釈できる。

この意味で、(F) は (E) と矛盾せず、

\begin{equation}
    \boxed{
        \text{Solver の構造的一意性}
        \longrightarrow
        \text{離散評価}
        \longrightarrow
        \text{細分化極限}
    }
\end{equation}

という流れを構成する。


%------------------------------------------------------------
\subsection{(F.13) 本節における証明事項と論考事項の区別}
\label{subsec:F13}

本節では、以下を明確に区別する。

\begin{enumerate}

    \item
    可測分割

    \begin{equation}
        \mathcal P_N
    \end{equation}

    を考え、その測度要素

    \begin{equation}
        \Delta\mu_j^{(N)}
        =
        \mu(E_j^{(N)})
    \end{equation}

    を定義すること。

    \item
    細分化メッシュ

    \begin{equation}
        \left\|
        \mathcal P_N
        \right\|_\mu
        =
        \max_j
        \Delta\mu_j^{(N)}
    \end{equation}

    を定義すること。

    \item
    メッシュが

    \begin{equation}
        \left\|
        \mathcal P_N
        \right\|_\mu
        \to0
    \end{equation}

    となること、および積分近似が収束することを仮定すれば、

    \begin{equation}
        T_N
        \to
        \int_I BIK\,d\mu
    \end{equation}

    が従うこと。

    \item
    一方で、任意の ECA Solver 集合について実際に

    \begin{equation}
        \left\|
        \mathcal P_N
        \right\|_\mu
        \to0
    \end{equation}

    となる細分化列を構成できることは、
    本節の定理から自動的には導かれないこと。

\end{enumerate}

したがって、

\begin{equation}
    \boxed{
        \text{離散極限定理}
    }
\end{equation}

と

\begin{equation}
    \boxed{
        \text{ECA Solver 集合における細分化極限の実現可能性}
    }
\end{equation}

を分離する。

後者については、本節では論考として扱う。


%------------------------------------------------------------
\subsection{(F.14) 本節の結論}
\label{subsec:F14}

以上により、ECA の離散形式

\begin{equation}
    T_N
    =
    \sum_{j=1}^{N}
    B_jI_jK_j
    \Delta\mu_j^{(N)}
\end{equation}

を、

\begin{equation}
    f(x)=B(x)I(x)K(x)
\end{equation}

によって

\begin{equation}
    T_N
    =
    \sum_{j=1}^{N}
    f\!\left(x_j^{(N)}\right)
    \Delta\mu_j^{(N)}
\end{equation}

と表現できる。

さらに、

\begin{equation}
    \left\|
    \mathcal P_N
    \right\|_\mu
    =
    \max_j\Delta\mu_j^{(N)}
    \longrightarrow0
\end{equation}

および対応する積分近似の収束を仮定すれば、

\begin{equation}
    \boxed{
        \lim_{N\to\infty}
        T_N
        =
        \int_I
        B(x)I(x)K(x)\,d\mu(x)
    }
\end{equation}

が成立する。

一様正規化の場合には、

\begin{equation}
    \Delta\mu_j^{(N)}
    =
    \frac1N
\end{equation}

であり、

\begin{equation}
    \left\|
    \mathcal P_N
    \right\|_\mu
    =
    \frac1N
    \longrightarrow0
\end{equation}

となるため、ECA の既存離散式と連続形式との対応が
直接的に確認できる。

したがって、本節では、

\begin{equation}
    \boxed{
        \text{離散 Solver の細分化}
        \longrightarrow
        \text{測度メッシュの消失}
        \longrightarrow
        \text{連続評価への極限}
    }
\end{equation}

という構造を確立する。

%################################################################
% (G)
%################################################################

\section{(G) Solver 集合における一般解の連続的構造}
\label{sec:G}

%------------------------------------------------------------
\subsection{(G.1) 本節の目的}
\label{subsec:G1}

(A)--(F) において、Solver 集合、評価空間、一般解、
評価写像、正規化測度および離散 Solver の細分化極限について
定義および定式化を行った。

(G) では、これらの構造を一般解 $\mathcal{G}$ のもとに統合し、
Solver 集合上の一般解が持つ連続的構造を記述する。

本節において重要なのは、(F) で扱った離散極限そのものを
新たに証明することではない。

(F) において、

\begin{equation}
    T_N
    =
    \sum_{j=1}^{N}
    f\!\left(x_j^{(N)}\right)
    \Delta\mu_j^{(N)}
\end{equation}

に対して適切な細分化条件および積分近似の収束条件を仮定した場合、

\begin{equation}
    \lim_{N\to\infty}T_N
    =
    \int_I f(x)\,d\mu(x)
\end{equation}

が成立することを定式化した。

したがって、本節ではこの結果を前提として、

\begin{equation}
    \mathcal{G}:\mathcal{S}\longrightarrow X
\end{equation}

という一般解の構造が、離散評価と連続評価の間をどのように
媒介するかを記述する。


%------------------------------------------------------------
\subsection{(G.2) (A)--(F) から引き継がれる構造}
\label{subsec:G2}

(A)--(D) において構成された Solver 集合を

\begin{equation}
    \mathcal{S}
\end{equation}

とする。

また、(B) において構成された評価空間を

\begin{equation}
    X
\end{equation}

とする。

(B)--(E) において扱った評価写像を

\begin{equation}
    \Phi:
    \mathcal{S}
    \longrightarrow
    X
\end{equation}

とする。

さらに、(D) において一般解を

\begin{equation}
    \mathcal{G}:
    \mathcal{S}
    \longrightarrow
    X
\end{equation}

として定義した。

(E) においては、Solver 同型

\begin{equation}
    T:
    \mathcal{S}_1
    \overset{\sim}{\longrightarrow}
    \mathcal{S}_2
\end{equation}

に対して、一般解が

\begin{equation}
    \mathcal{G}_2\circ T
    =
    \mathcal{G}_1
\end{equation}

という可換性を満たすことを一般解の構造的一意性として扱った。

一方、(F) においては、有限段階の可測分割

\begin{equation}
    \mathcal{P}_N
    =
    \left\{
        E_1^{(N)},
        \ldots,
        E_N^{(N)}
    \right\}
\end{equation}

および測度要素

\begin{equation}
    \Delta\mu_j^{(N)}
    =
    \mu\!\left(E_j^{(N)}\right)
\end{equation}

を用いて、離散評価を

\begin{equation}
    T_N
    =
    \sum_{j=1}^{N}
    f\!\left(x_j^{(N)}\right)
    \Delta\mu_j^{(N)}
\end{equation}

と記述した。

ここで、

\begin{equation}
    f(x)
    =
    B(x)I(x)K(x)
\end{equation}

である。

これらを基礎として、以下では一般解 $\mathcal{G}$ と
連続評価との関係を定式化する。


%------------------------------------------------------------
\subsection{(G.3) 一般解 $\mathcal{G}$ の定義}
\label{subsec:G3}

一般解 $\mathcal{G}$ は、Solver に対して最終的な数値を
直接割り当てる写像としてではなく、Solver が持つ評価構造を
評価空間 $X$ の元として表現する写像として定義する。

したがって、

\begin{equation}
    \boxed{
        \mathcal{G}:
        \mathcal{S}
        \longrightarrow
        X
    }
    \label{eq:G-general-solution}
\end{equation}

とする。

任意の

\begin{equation}
    S\in\mathcal{S}
\end{equation}

に対して、

\begin{equation}
    \mathcal{G}(S)\in X
\end{equation}

である。

したがって、

\begin{equation}
    \boxed{
        S
        \longmapsto
        \mathcal{G}(S)\in X
    }
\end{equation}

が一般解の基本的な写像構造となる。

ここでいう一般解は、特定の Solver における個別の解を
列挙するものではない。

一般解 $\mathcal{G}$ は、Solver 集合の各元が持つ評価構造を
評価空間上に表現することによって、Solver 集合全体に共通する
一般的な評価構造を記述するものである。


%------------------------------------------------------------
\subsection{(G.4) 評価構造としての一般解}
\label{subsec:G4}

ECA において評価構造が

\begin{equation}
    B,\qquad I,\qquad K
\end{equation}

によって表現される場合を考える。

評価空間 $X$ が

\begin{equation}
    X
    =
    X_B\times X_I\times X_K
\end{equation}

として構成できる場合には、

\begin{equation}
    \boxed{
        \mathcal{G}(S)
        =
        \left(
            B(S),
            I(S),
            K(S)
        \right)
    }
    \label{eq:G-BIK}
\end{equation}

と表現できる。

ただし、この表現は評価空間が実際に
$B,I,K$ の直積として構成される場合の具体的表現である。

したがって、

\begin{equation}
    \mathcal{G}:
    \mathcal{S}\longrightarrow X
\end{equation}

という一般解の定義そのものを、

\begin{equation}
    \mathcal{G}(S)
    =
    \left(B(S),I(S),K(S)\right)
\end{equation}

という特定の表現に限定するものではない。

一般解の本質は、

\begin{equation}
    \boxed{
        S\longmapsto\mathcal{G}(S)\in X
    }
\end{equation}

という写像構造にある。


%------------------------------------------------------------
\subsection{(G.5) 一般解から最終評価への写像}
\label{subsec:G5}

一般解 $\mathcal{G}(S)$ は評価空間 $X$ の元であるため、
そこから最終的なスカラー評価を取り出す写像を

\begin{equation}
    \Psi:
    X
    \longrightarrow
    \mathbb{R}
\end{equation}

とする。

例えば、

\begin{equation}
    x=(B,I,K)
\end{equation}

として表現される場合には、

\begin{equation}
    \Psi(B,I,K)
    =
    BIK
\end{equation}

と定義できる。

このとき、Solver $S$ に対する最終評価を

\begin{equation}
    g(S)
    :=
    \Psi\!\left(\mathcal{G}(S)\right)
\end{equation}

と定義する。

すなわち、

\begin{equation}
    \boxed{
        g
        =
        \Psi\circ\mathcal{G}
    }
    \label{eq:G-composed-map}
\end{equation}

である。

したがって、

\begin{equation}
    g:
    \mathcal{S}
    \longrightarrow
    \mathbb{R}
\end{equation}

となり、

\begin{equation}
    \boxed{
        g(S)
        =
        \Psi\!\left(\mathcal{G}(S)\right)
    }
    \label{eq:G-scalar-evaluation}
\end{equation}

を得る。

以上から、一般解と最終評価の関係は

\begin{equation}
    \boxed{
        \mathcal{S}
        \xrightarrow{\mathcal{G}}
        X
        \xrightarrow{\Psi}
        \mathbb{R}
    }
    \label{eq:G-evaluation-chain}
\end{equation}

という二段階の写像として記述される。


%------------------------------------------------------------
\subsection{(G.6) 評価写像 $\Phi$ と一般解 $\mathcal{G}$ の関係}
\label{subsec:G6}

(B)--(E) において導入された評価写像を

\begin{equation}
    \Phi:
    \mathcal{S}
    \longrightarrow
    X
\end{equation}

とする。

一方、一般解は

\begin{equation}
    \mathcal{G}:
    \mathcal{S}
    \longrightarrow
    X
\end{equation}

である。

したがって、両者は型として同じ写像を持つ。

しかし、

\begin{equation}
    \Phi=\mathcal{G}
\end{equation}

であることを一般解の定義そのものとして要求する必要はない。

$\Phi$ は Solver を評価空間へ対応させる評価写像であり、

\begin{equation}
    \mathcal{G}(S)
\end{equation}

は Solver $S$ の評価構造を一般解として表現するものである。

両者が同一の評価構造を表現する場合には、

\begin{equation}
    \mathcal{G}(S)
    =
    \Phi(S)
    \qquad
    (\forall S\in\mathcal{S})
\end{equation}

として同一視できる。

この場合には、

\begin{equation}
    \mathcal{G}
    =
    \Phi
\end{equation}

となる。

したがって、(G) において必要なのは
$\mathcal{G}$ と $\Phi$ の名称を機械的に同一視することではなく、

\begin{equation}
    \boxed{
        \mathcal{G}:
        \mathcal{S}\to X
    }
\end{equation}

という一般解の構造を明確にすることである。


%------------------------------------------------------------
\subsection{(G.7) Solver 同型性と一般解}
\label{subsec:G7}

(E) において定義された Solver 同型

\begin{equation}
    T:
    \mathcal{S}_1
    \overset{\sim}{\longrightarrow}
    \mathcal{S}_2
\end{equation}

を考える。

それぞれの Solver 集合に対する一般解を

\begin{equation}
    \mathcal{G}_1:
    \mathcal{S}_1\to X_1,
    \qquad
    \mathcal{G}_2:
    \mathcal{S}_2\to X_2
\end{equation}

とする。

(E) における強い一意性の条件に従い、
評価空間の対応を同一視できる場合には、

\begin{equation}
    \boxed{
        \mathcal{G}_2\circ T
        =
        \mathcal{G}_1
    }
    \label{eq:G-isomorphism}
\end{equation}

が成立する。

すなわち、任意の

\begin{equation}
    S\in\mathcal{S}_1
\end{equation}

について、

\begin{equation}
    \mathcal{G}_2(T(S))
    =
    \mathcal{G}_1(S)
\end{equation}

である。

この式は、一般解が Solver の具体的な表現そのものに
依存するのではなく、(E) において定められた Solver 構造を
保存する対応によって定まることを表す。


%------------------------------------------------------------
\subsection{(G.8) 最終評価の不変性}
\label{subsec:G8}

一般解に対応する最終評価を

\begin{equation}
    g_1
    =
    \Psi_1\circ\mathcal{G}_1,
    \qquad
    g_2
    =
    \Psi_2\circ\mathcal{G}_2
\end{equation}

とする。

評価空間間の対応

\begin{equation}
    U:
    X_1
    \longrightarrow
    X_2
\end{equation}

が存在し、

\begin{equation}
    U\circ\mathcal{G}_1
    =
    \mathcal{G}_2\circ T
\end{equation}

および

\begin{equation}
    \Psi_2\circ U
    =
    \Psi_1
\end{equation}

を満たすとする。

このとき、

\begin{align}
    g_2\circ T
    &=
    \Psi_2\circ\mathcal{G}_2\circ T
    \\
    &=
    \Psi_2\circ U\circ\mathcal{G}_1
    \\
    &=
    \Psi_1\circ\mathcal{G}_1
    \\
    &=
    g_1
\end{align}

である。

したがって、

\begin{equation}
    \boxed{
        g_2\circ T
        =
        g_1
    }
    \label{eq:G-scalar-invariance}
\end{equation}

を得る。

これは、Solver の表現を変更しても、
Solver 構造および最終評価構造が保存される限り、
最終評価が対応して保存されることを意味する。


%------------------------------------------------------------
\subsection{(G.9) 離散 Solver における一般解}
\label{subsec:G9}

(F) において定義した可測分割を

\begin{equation}
    \mathcal{P}_N
    =
    \left\{
        E_1^{(N)},
        \ldots,
        E_N^{(N)}
    \right\}
\end{equation}

とする。

各分割要素 $E_j^{(N)}$ に対応する Solver を

\begin{equation}
    S_j^{(N)}
\end{equation}

とし、その一般解を

\begin{equation}
    x_j^{(N)}
    :=
    \mathcal{G}\!\left(S_j^{(N)}\right)
    \in X
\end{equation}

と定義する。

このとき、最終評価は

\begin{equation}
    \Psi\!\left(
        x_j^{(N)}
    \right)
    =
    \Psi\!\left(
        \mathcal{G}\!\left(S_j^{(N)}\right)
    \right)
\end{equation}

である。

したがって、(F) における離散一般解は

\begin{equation}
    \boxed{
        T_N
        =
        \sum_{j=1}^{N}
        \Psi\!\left(
            \mathcal{G}\!\left(S_j^{(N)}\right)
        \right)
        \Delta\mu_j^{(N)}
    }
    \label{eq:G-discrete-general-solution}
\end{equation}

と表現できる。

ここで、

\begin{equation}
    g
    =
    \Psi\circ\mathcal{G}
\end{equation}

であるから、

\begin{equation}
    T_N
    =
    \sum_{j=1}^{N}
    g\!\left(S_j^{(N)}\right)
    \Delta\mu_j^{(N)}
\end{equation}

とも書ける。


%------------------------------------------------------------
\subsection{(G.10) 離散一般解から連続一般解への接続}
\label{subsec:G10}

(F) の細分化極限定理を適用する。

すなわち、

\begin{equation}
    \left\|
    \mathcal{P}_N
    \right\|_\mu
    \longrightarrow0
\end{equation}

および対応する積分近似の収束条件が成立すると仮定する。

このとき、

\begin{equation}
    \lim_{N\to\infty}
    T_N
    =
    \int_I
    f(x)\,d\mu(x)
\end{equation}

である。

ここで、一般解による評価構造を用いて

\begin{equation}
    f(x)
    =
    \Psi\!\left(\mathcal{G}(x)\right)
\end{equation}

と表現できる場合には、

\begin{equation}
    \boxed{
        \lim_{N\to\infty}
        T_N
        =
        \int_I
        \Psi\!\left(\mathcal{G}(x)\right)
        \,d\mu(x)
    }
    \label{eq:G-continuous-limit}
\end{equation}

となる。

したがって、離散 Solver に対する一般解の評価は、
細分化極限を通じて連続的な評価へ接続される。

この関係を記号的に

\begin{equation}
    \boxed{
        \mathcal{G}_{\mathrm{discrete}}
        \longrightarrow
        \mathcal{G}_{\mathrm{continuous}}
    }
\end{equation}

と表すことができる。


%------------------------------------------------------------
\subsection{(G.11) 一般解の連続的構造}
\label{subsec:G11}

(G.10) により、一般解 $\mathcal{G}$ による評価は
離散的な Solver の列に限定されず、適切な細分化極限のもとで
連続的な評価空間上に拡張される。

すなわち、離散段階における

\begin{equation}
    S_j^{(N)}
    \longmapsto
    \mathcal{G}\!\left(S_j^{(N)}\right)
\end{equation}

という対応に対して、細分化極限では

\begin{equation}
    x
    \longmapsto
    \mathcal{G}(x)
\end{equation}

という連続的評価構造を考えることができる。

ここで重要なのは、(F) で述べたように、

\begin{equation}
    \left\|
    \mathcal{P}_N
    \right\|_\mu
    \longrightarrow0
\end{equation}

という条件そのものを本節で新たに証明しているわけではないことである。

(G) においては、(F) の極限結果を前提として、
その極限が一般解 $\mathcal{G}$ による評価構造を
連続的評価へ接続することを記述している。


%------------------------------------------------------------
\subsection{(G.12) 一般解の数学的意味}
\label{subsec:G12}

以上の構成から、一般解 $\mathcal{G}$ は単に Solver に
数値を割り当てる関数ではない。

一般解は、

\begin{equation}
    \boxed{
        \mathcal{G}:
        \mathcal{S}\longrightarrow X
    }
\end{equation}

によって Solver 集合上の評価構造を評価空間上へ写像する。

さらに、

\begin{equation}
    \boxed{
        \Psi:
        X\longrightarrow\mathbb{R}
    }
\end{equation}

によって、その評価構造から最終的なスカラー評価を取り出す。

したがって、

\begin{equation}
    \boxed{
        g
        =
        \Psi\circ\mathcal{G}
    }
\end{equation}

である。

この構造によって、

\begin{equation}
    \boxed{
        \mathcal{S}
        \xrightarrow{\mathcal{G}}
        X
        \xrightarrow{\Psi}
        \mathbb{R}
    }
\end{equation}

という統一された写像構造が得られる。


%------------------------------------------------------------
\subsection{(G.13) (E) および (F) との統合}
\label{subsec:G13}

(E) においては、Solver 同型に対して一般解が

\begin{equation}
    \mathcal{G}_2\circ T
    =
    \mathcal{G}_1
\end{equation}

という可換性を持つことにより、
Solver の表現に依存しない一般解の構造的一意性を扱った。

(F) においては、離散 Solver の評価を

\begin{equation}
    T_N
    =
    \sum_{j=1}^{N}
    \Psi\!\left(
        \mathcal{G}\!\left(S_j^{(N)}\right)
    \right)
    \Delta\mu_j^{(N)}
\end{equation}

と記述し、適切な細分化極限のもとで

\begin{equation}
    T_N
    \longrightarrow
    \int_I
    \Psi\!\left(\mathcal{G}(x)\right)
    \,d\mu(x)
\end{equation}

という連続形式への接続を定式化した。

したがって、(G) においてこれらを統合すると、

\begin{equation}
    \boxed{
        \text{Solver 構造}
        \longrightarrow
        \text{一般解 }\mathcal{G}
        \longrightarrow
        \text{離散評価}
        \longrightarrow
        \text{連続評価}
    }
\end{equation}

という一連の構造が得られる。


%------------------------------------------------------------
\subsection{(G.14) (A)--(F) との関係}
\label{subsec:G14}

(A) では Solver 集合の数学的構造を定義した。

(B) では、その Solver 集合を評価するための一般評価空間
$X$ および評価写像を構成した。

(C) では、最終評価関数に必要となる測度論的性質を扱った。

(D) では、

\begin{equation}
    \mathcal{G}:
    \mathcal{S}\longrightarrow X
\end{equation}

という一般解の存在を定式化した。

(E) では、Solver 同型によって対応する Solver に対して
一般解が可換することを定式化し、一般解の構造的一意性を扱った。

(F) では、Solver の離散的評価を可測分割と正規化測度によって
記述し、適切な細分化極限のもとで連続評価へ接続した。

以上を踏まえると、(G) における一般解は、

\begin{equation}
    \boxed{
        \mathcal{G}:
        \mathcal{S}\longrightarrow X
    }
\end{equation}

という基本構造を保持しながら、

\begin{equation}
    \boxed{
        \mathcal{S}
        \xrightarrow{\mathcal{G}}
        X
        \xrightarrow{\Psi}
        \mathbb{R}
    }
\end{equation}

という評価構造を形成し、さらに (F) の細分化極限によって
連続的評価へ接続される。


%------------------------------------------------------------
\subsection{(G.15) (H) への接続}
\label{subsec:G15}

(G) において一般解

\begin{equation}
    \mathcal{G}:
    \mathcal{S}\longrightarrow X
\end{equation}

の構造が定まることにより、Solver の内部構造と一般解との
対応をさらに具体化することが可能となる。

特に、Solver が持つ公理選択、離散的構成および連続的構成を
一般解の像からどの程度復元できるかという問題を考えることができる。

この問題は、(H) において扱う Solver の構造的情報と
一般解との対応へ接続される。

したがって、(G) から (H) への構造は、

\begin{equation}
    \boxed{
        \text{Solver}
        \xrightarrow{\mathcal{G}}
        \text{一般評価構造}
        \longrightarrow
        \text{Solver 構造の復元}
    }
\end{equation}

として理解される。

%################################################################
% (H)
%################################################################

\section{(H) 中間 Solver の厳密な定義と構造}
\label{sec:H}

%------------------------------------------------------------
\subsection{(H.1) 基本設定}
\label{subsec:H1}

(A)--(G) において定義された Solver 集合、一般解写像、
評価空間および ECA の評価構造を前提とする。

Solver 集合を

\begin{equation}
    \mathcal S
    =
    \left\{
        s
        \mid
        s=(\mathcal A_s,\mathcal R_s)
    \right\}
    \label{eq:H-solver}
\end{equation}

とする。

ここで、

\begin{equation}
    \mathcal A_s
\end{equation}

は Solver $s$ において採用される公理の選択集合、

\begin{equation}
    \mathcal R_s
\end{equation}

は推論規則および構成制約の集合を表す。

また、Solver $s$ によって表現される理論を与える写像を

\begin{equation}
    \mathcal T:
    \mathcal S
    \longrightarrow
    \mathrm{Th}
    \label{eq:H-theory-map}
\end{equation}

とする。

したがって、

\begin{equation}
    \mathcal T(s)
\end{equation}

は公理集合そのものではなく、Solver $s$ の公理選択および
推論規則・構成制約によって表現される理論を表す。

一方、(G) において定義された一般解写像を

\begin{equation}
    \boxed{
        \mathcal G:
        \mathcal S
        \longrightarrow
        X
    }
    \label{eq:H-general-solution}
\end{equation}

とする。


%------------------------------------------------------------
\subsection{(H.2) ECA の公理選択空間}
\label{subsec:H2}

ECA における公理選択空間を

\begin{equation}
    I
    \subseteq
    \mathcal P(S)
    \label{eq:H-axiom-space}
\end{equation}

とする。

その元

\begin{equation}
    i\in I
\end{equation}

を、ECA において実際に採用され得る公理集合とする。

さらに、ECA において定義された離散的要素集合を

\begin{equation}
    B(i)
    \subseteq
    i
    \subseteq
    S
    \label{eq:H-discrete-subset}
\end{equation}

とする。

ここで $B(i)$ は、公理選択 $i$ に対応して取り出される
離散的評価対象を表す。


%------------------------------------------------------------
\subsection{(H.3) Solver--ECA 適合条件}
\label{subsec:H3}

Solver の公理選択成分と ECA の公理選択空間を接続するため、
次の条件を定義する。

\begin{definition}[Solver--ECA 適合条件]
\label{def:H-compatible}
Solver $s\in\mathcal S$ が ECA の公理選択空間 $I$ に適合するとは、

\begin{equation}
    \boxed{
        \mathcal A_s\in I
    }
    \label{eq:H-compatible}
\end{equation}

が成立することをいう。
\end{definition}

したがって、ECA に適合する Solver 全体を

\begin{equation}
    \boxed{
        \mathcal S_I
        :=
        \left\{
            s\in\mathcal S
            \mid
            \mathcal A_s\in I
        \right\}
    }
    \label{eq:H-compatible-solvers}
\end{equation}

と定義する。


%------------------------------------------------------------
\subsection{(H.4) 公理選択復元写像}
\label{subsec:H4}

\begin{definition}[公理選択復元写像]
\label{def:H-axiom-map}

ECA に適合する Solver

\begin{equation}
    s\in\mathcal S_I
\end{equation}

に対して、

\begin{equation}
    \mathfrak A:
    \mathcal S_I
    \longrightarrow
    I
\end{equation}

を

\begin{equation}
    \boxed{
        \mathfrak A(s)
        :=
        \mathcal A_s
    }
    \label{eq:H-axiom-reconstruction}
\end{equation}

によって定義する。

以下、

\begin{equation}
    \boxed{
        i_s
        :=
        \mathfrak A(s)
        =
        \mathcal A_s
    }
    \label{eq:H-is}
\end{equation}

と記す。
\end{definition}

この定義により、公理選択 $i_s$ は外部から恣意的に
与えられるものではなく、Solver

\begin{equation}
    s=(\mathcal A_s,\mathcal R_s)
\end{equation}

の内部構造に含まれる公理選択成分から直接得られる。


%------------------------------------------------------------
\subsection{(H.5) 公理選択復元写像の well-defined 性}
\label{subsec:H5}

\begin{proposition}[公理選択復元写像の well-defined 性]
\label{prop:H-axiom-well-defined}

写像

\begin{equation}
    \mathfrak A:
    \mathcal S_I
    \longrightarrow
    I
\end{equation}

は well-defined である。
\end{proposition}

\begin{proof}

任意の

\begin{equation}
    s\in\mathcal S_I
\end{equation}

を取る。

$\mathcal S_I$ の定義より、

\begin{equation}
    \mathcal A_s\in I
\end{equation}

が成立する。

したがって、

\begin{equation}
    \mathfrak A(s)
    :=
    \mathcal A_s
\end{equation}

は $I$ の元を与える。

また、Solver は

\begin{equation}
    s=(\mathcal A_s,\mathcal R_s)
\end{equation}

として定義されているため、同一の Solver $s$ に対する
$\mathcal A_s$ はその定義によって一意に定まる。

よって、

\begin{equation}
    \mathfrak A(s)=\mathcal A_s
\end{equation}

は各 $s\in\mathcal S_I$ に対して一意に定まり、

\begin{equation}
    \mathfrak A:
    \mathcal S_I\to I
\end{equation}

は well-defined である。

\end{proof}


%------------------------------------------------------------
\subsection{(H.6) $\mathcal T(s)$ と公理選択の関係}
\label{subsec:H6}

Solver

\begin{equation}
    s=(\mathcal A_s,\mathcal R_s)
\end{equation}

に対して、

\begin{equation}
    \mathcal T(s)
    =
    \mathcal T(\mathcal A_s,\mathcal R_s)
    \label{eq:H-theory-decomposition}
\end{equation}

と表記する。

ここで重要なのは、

\begin{equation}
    \mathcal T(s)
\end{equation}

から公理集合を単純な集合論的逆像として復元する必要はない
ということである。

Solver の定義そのものが

\begin{equation}
    s=(\mathcal A_s,\mathcal R_s)
\end{equation}

を与えているため、公理選択成分は

\begin{equation}
    \mathcal A_s
\end{equation}

として直接抽出できる。

したがって、ECA に適合する Solver に対して、

\begin{equation}
    \boxed{
        i_s
        =
        \mathfrak A(s)
        =
        \mathcal A_s
    }
    \label{eq:H-theory-axiom}
\end{equation}

である。

従って、

\begin{equation}
    \boxed{
        \mathcal T(s)
        =
        \mathcal T(i_s,\mathcal R_s)
    }
    \label{eq:H-theory-is}
\end{equation}

と表現できる。

ここで $\mathcal T(s)$ は理論を表し、
$i_s$ はその Solver における公理選択を表す。

両者を同一視しないことが重要である。


%------------------------------------------------------------
\subsection{(H.7) Solver 同値性と公理選択}
\label{subsec:H7}

Solver 集合上の同値関係を

\begin{equation}
    s\sim_{\mathrm{Sol}}t
\end{equation}

とする。

Solver として同値であっても、その内部表現が異なることは
許容される。

したがって、

\begin{equation}
    s\sim_{\mathrm{Sol}}t
\end{equation}

から直ちに

\begin{equation}
    \mathcal A_s
    =
    \mathcal A_t
\end{equation}

を要求する必要はない。

そこで、公理選択空間 $I$ 上に同値関係

\begin{equation}
    \sim_I
\end{equation}

を導入し、Solver 同値性と公理選択との間に

\begin{equation}
    \boxed{
        s\sim_{\mathrm{Sol}}t
        \Longrightarrow
        \mathfrak A(s)\sim_I\mathfrak A(t)
    }
    \label{eq:H-solver-axiom-equivalence}
\end{equation}

という整合条件を置く。

これは、Solver の内部表現が異なっていても、
ECA における公理選択の評価構造が保存されることを意味する。


%------------------------------------------------------------
\subsection{(H.8) Solver の離散評価構造}
\label{subsec:H8}

任意の

\begin{equation}
    i\in I
\end{equation}

に対して、

\begin{equation}
    B(i)\subseteq i\subseteq S
\end{equation}

が成立する。

したがって、Solver $s\in\mathcal S_I$ に対応する公理選択

\begin{equation}
    i_s=\mathfrak A(s)
\end{equation}

について、

\begin{equation}
    B(i_s)
    \subseteq
    i_s
    \subseteq
    S
\end{equation}

が成立する。

\begin{definition}[Solver の離散評価構造]
\label{def:H-discrete-structure}

Solver $s\in\mathcal S_I$ に対して、その離散評価構造を

\begin{equation}
    \boxed{
        D_s
        :=
        B(i_s)
        =
        B(\mathfrak A(s))
    }
    \label{eq:H-Ds}
\end{equation}

と定義する。
\end{definition}

したがって、

\begin{equation}
    \boxed{
        D_s\subseteq i_s\subseteq S
    }
    \label{eq:H-Ds-inclusion}
\end{equation}

が直ちに成立する。

$D_s$ が有限または可算であり、
離散評価に用いる重み関数

\begin{equation}
    w_s:
    S\longrightarrow[0,1]
\end{equation}

が

\begin{equation}
    \sum_{u\in S}w_s(u)=1
    \label{eq:H-weight-normalization}
\end{equation}

を満たす場合には、Solver $s$ の離散評価を

\begin{equation}
    \boxed{
        E_{\mathrm d}(s)
        :=
        \sum_{b\in D_s}w_s(b)
    }
    \label{eq:H-discrete-evaluation}
\end{equation}

と定義する。


%------------------------------------------------------------
\subsection{(H.9) 離散評価の基本性質}
\label{subsec:H9}

\begin{proposition}[離散評価の有界性]
\label{prop:H-discrete-bound}

$D_s$ が有限または可算であり、

\begin{equation}
    w_s:S\to[0,1],
    \qquad
    \sum_{u\in S}w_s(u)=1
\end{equation}

を満たすならば、

\begin{equation}
    \boxed{
        0
        \leq
        E_{\mathrm d}(s)
        \leq
        1
    }
    \label{eq:H-discrete-bound}
\end{equation}

である。
\end{proposition}

\begin{proof}

各 $b\in D_s$ に対して

\begin{equation}
    w_s(b)\geq0
\end{equation}

であるから、

\begin{equation}
    E_{\mathrm d}(s)
    =
    \sum_{b\in D_s}w_s(b)
    \geq0.
\end{equation}

また、

\begin{equation}
    D_s\subseteq S
\end{equation}

であるため、

\begin{equation}
    \sum_{b\in D_s}w_s(b)
    \leq
    \sum_{u\in S}w_s(u)
    =
    1.
\end{equation}

したがって、

\begin{equation}
    0\leq E_{\mathrm d}(s)\leq1.
\end{equation}

\end{proof}


%------------------------------------------------------------
\subsection{(H.10) Solver の連続評価構造：測度そのもの}
\label{subsec:H10}

Solver $s$ に対応する公理選択 $i_s$ 上で連続的評価を行うため、
測度

\begin{equation}
    \mu_s
\end{equation}

を考える。

\begin{definition}[連続評価測度]
\label{def:H-continuous-measure}

Solver $s$ に対して、連続評価に用いられる測度そのものを

\begin{equation}
    \boxed{
        C_s^{(\mu)}
        :=
        \mu_s
    }
    \label{eq:H-Cs-measure}
\end{equation}

と定義する。

対応する連続評価値を

\begin{equation}
    \boxed{
        E_{\mathrm c}(s)
        :=
        \mu_s(i_s)
        =
        \int_S
        \mathbf 1_{i_s}(u)
        \,d\mu_s(u)
    }
    \label{eq:H-continuous-evaluation}
\end{equation}

とする。

ここで

\begin{equation}
    \mathbf1_{i_s}:S\to\{0,1\}
\end{equation}

は $i_s$ の指示関数である。
\end{definition}


%------------------------------------------------------------
\subsection{(H.11) Solver の連続評価構造：測度空間}
\label{subsec:H11}

測度そのものだけでは、可測性や測度の定義域を必要とする
場合に十分でない。

したがって、可測空間および測度を含めて連続評価構造を扱う場合には、
次の定義を用いる。

\begin{definition}[連続評価測度空間]
\label{def:H-continuous-measure-space}

Solver $s$ に対して、連続評価を成立させる測度空間を

\begin{equation}
    \boxed{
        C_s^{(\mathrm{ms})}
        :=
        (I_s,\Sigma_{I_s},\mu_s)
    }
    \label{eq:H-Cs-measure-space}
\end{equation}

と定義する。

ここで、

\begin{equation}
    I_s
\end{equation}

は連続評価に用いられる公理選択空間、

\begin{equation}
    \Sigma_{I_s}
\end{equation}

は $I_s$ 上の $\sigma$-代数、

\begin{equation}
    \mu_s:
    \Sigma_{I_s}
    \longrightarrow
    [0,1]
\end{equation}

は測度である。
\end{definition}

したがって、

\begin{equation}
    C_s^{(\mu)}
\end{equation}

と

\begin{equation}
    C_s^{(\mathrm{ms})}
\end{equation}

は異なるレベルの構造を表す。

すなわち、

\begin{equation}
    C_s^{(\mu)}
\end{equation}

は評価作用そのものを表し、

\begin{equation}
    C_s^{(\mathrm{ms})}
\end{equation}

はその評価作用を成立させる測度論的構造全体を表す。


%------------------------------------------------------------
\subsection{(H.12) 測度と測度空間の関係}
\label{subsec:H12}

測度空間

\begin{equation}
    C_s^{(\mathrm{ms})}
    =
    (I_s,\Sigma_{I_s},\mu_s)
\end{equation}

から測度成分を取り出す写像

\begin{equation}
    \pi_\mu:
    C_s^{(\mathrm{ms})}
    \longrightarrow
    C_s^{(\mu)}
\end{equation}

を

\begin{equation}
    \boxed{
        \pi_\mu
        (I_s,\Sigma_{I_s},\mu_s)
        =
        \mu_s
    }
    \label{eq:H-measure-projection}
\end{equation}

によって定義する。

したがって、

\begin{equation}
    C_s^{(\mathrm{ms})}
    \longrightarrow
    C_s^{(\mu)}
\end{equation}

という自然な忘却写像が存在する。

一方、一般には

\begin{equation}
    \mu_s
\end{equation}

のみから

\begin{equation}
    (I_s,\Sigma_{I_s})
\end{equation}

を一意に復元することはできない。

したがって、

\begin{equation}
    \boxed{
        C_s^{(\mu)}
        \text{ は }
        C_s^{(\mathrm{ms})}
        \text{ の情報を完全には保持しない}
    }
    \label{eq:H-measure-information-loss}
\end{equation}

。


%------------------------------------------------------------
\subsection{(H.13) 離散評価と連続評価の統一}
\label{subsec:H13}

Solver $s\in\mathcal S_I$ に対して、

\begin{equation}
    i_s=\mathfrak A(s)
\end{equation}

と置く。

このとき、離散評価は

\begin{equation}
    E_{\mathrm d}(s)
    =
    \sum_{b\in B(i_s)}w_s(b)
\end{equation}

であり、連続評価は

\begin{equation}
    E_{\mathrm c}(s)
    =
    \int_S
    \mathbf1_{i_s}(u)\,d\mu_s(u)
\end{equation}

である。

したがって、

\begin{equation}
    \boxed{
    \begin{aligned}
        i_s
        &=
        \mathfrak A(s),
        \\[1mm]
        D_s
        &=
        B(i_s),
        \\[1mm]
        E_{\mathrm d}(s)
        &=
        \sum_{b\in D_s}w_s(b),
        \\[1mm]
        C_s^{(\mu)}
        &=
        \mu_s,
        \\[1mm]
        E_{\mathrm c}(s)
        &=
        \int_S
        \mathbf1_{i_s}(u)\,d\mu_s(u).
    \end{aligned}
    }
    \label{eq:H-unified-evaluation}
\end{equation}

という統一的な構成が得られる。

ここで重要なのは、離散評価と連続評価が新しい二つの評価原理として
別々に導入されるわけではないことである。

同一の公理選択

\begin{equation}
    i_s
\end{equation}

に対して、ECA に既に定義された離散表現および連続表現を
適用することによって、両者が得られる。


%------------------------------------------------------------
\subsection{(H.14) 一般解における二重評価構造}
\label{subsec:H14}

(G) において定義された一般解

\begin{equation}
    \mathcal G:
    \mathcal S
    \longrightarrow
    X
\end{equation}

を考える。

Solver $s\in\mathcal S_I$ に対して、離散評価および連続評価を
一般解へ入力する構成写像

\begin{equation}
    \Phi_s:
    [0,1]\times[0,1]
    \longrightarrow
    X
\end{equation}

が存在し、

\begin{equation}
    \boxed{
        \mathcal G(s)
        =
        \Phi_s
        \left(
            E_{\mathrm d}(s),
            E_{\mathrm c}(s)
        \right)
    }
    \label{eq:H-general-solution-two-evaluations}
\end{equation}

と表現できる場合を考える。

このとき、一般解における離散評価および連続評価への依存性を
次のように定義する。


%------------------------------------------------------------
\subsubsection{(H.14.1) 離散評価への非自明な依存性}
\label{subsubsec:H14-1}

$\Phi_s$ が第1変数に非自明に依存するとは、

\begin{equation}
    \exists x_1,x_2,y\in[0,1]
\end{equation}

について

\begin{equation}
    x_1\neq x_2
\end{equation}

かつ

\begin{equation}
    \boxed{
        \Phi_s(x_1,y)
        \neq
        \Phi_s(x_2,y)
    }
    \label{eq:H-nontrivial-discrete}
\end{equation}

が成立することをいう。


%------------------------------------------------------------
\subsubsection{(H.14.2) 連続評価への非自明な依存性}
\label{subsubsec:H14-2}

$\Phi_s$ が第2変数に非自明に依存するとは、

\begin{equation}
    \exists x,y_1,y_2\in[0,1]
\end{equation}

について

\begin{equation}
    y_1\neq y_2
\end{equation}

かつ

\begin{equation}
    \boxed{
        \Phi_s(x,y_1)
        \neq
        \Phi_s(x,y_2)
    }
    \label{eq:H-nontrivial-continuous}
\end{equation}

が成立することをいう。


%------------------------------------------------------------
\subsection{(H.15) 中間 Solver の厳密な定義}
\label{subsec:H15}

\begin{definition}[中間 Solver]
\label{def:H-middle-solver}

Solver

\begin{equation}
    s\in\mathcal S_I
\end{equation}

が中間 Solver であるとは、以下の条件を満たすことをいう。

\begin{enumerate}[label=\textup{(\arabic*)}]

    \item
    ECA に適合する公理選択

    \begin{equation}
        i_s
        =
        \mathfrak A(s)
        =
        \mathcal A_s
        \in I
    \end{equation}

    が存在する。

    \item
    離散評価構造

    \begin{equation}
        D_s=B(i_s)
    \end{equation}

    が存在し、離散評価

    \begin{equation}
        E_{\mathrm d}(s)
        =
        \sum_{b\in D_s}w_s(b)
    \end{equation}

    が定義される。

    \item
    連続評価構造として、少なくとも

    \begin{equation}
        C_s^{(\mu)}
        =
        \mu_s
    \end{equation}

    が存在する。

    必要に応じて、より完全な構造として

    \begin{equation}
        C_s^{(\mathrm{ms})}
        =
        (I_s,\Sigma_{I_s},\mu_s)
    \end{equation}

    を用いる。

    \item
    一般解が

    \begin{equation}
        \mathcal G(s)
        =
        \Phi_s
        \left(
            E_{\mathrm d}(s),
            E_{\mathrm c}(s)
        \right)
    \end{equation}

    と表現され、$\Phi_s$ が第1変数および第2変数の双方に
    非自明に依存する。

\end{enumerate}
\end{definition}

この定義において重要なのは、中間 Solver を

\begin{equation}
    \text{離散構造}\cup\text{連続構造}
\end{equation}

として定義していないことである。

必要なのは、**同一の Solver に対して両方の評価構造が存在し、
しかもその双方が一般解の決定に実質的に寄与すること**である。


%------------------------------------------------------------
\subsection{(H.16) 中間 Solver 構造定理}
\label{subsec:H16}

\begin{theorem}[中間 Solver 構造定理]
\label{thm:H-middle-solver}

$s\in\mathcal S_I$ とする。

さらに、

\begin{equation}
    \mathcal G(s)
    =
    \Phi_s
    \left(
        E_{\mathrm d}(s),
        E_{\mathrm c}(s)
    \right)
\end{equation}

と表現できると仮定する。

このとき、$s$ が中間 Solver であることと、

\begin{equation}
    \exists x_1\neq x_2,\ y
\end{equation}

について

\begin{equation}
    \Phi_s(x_1,y)
    \neq
    \Phi_s(x_2,y)
    \label{eq:H-theorem-discrete}
\end{equation}

が成立し、かつ

\begin{equation}
    \exists x,\ y_1\neq y_2
\end{equation}

について

\begin{equation}
    \Phi_s(x,y_1)
    \neq
    \Phi_s(x,y_2)
    \label{eq:H-theorem-continuous}
\end{equation}

が成立することは同値である。
\end{theorem}

\begin{proof}

定義より、中間 Solver では一般解が離散評価と連続評価の
双方に非自明に依存する。

まず、離散評価について考える。

もし任意の

\begin{equation}
    x_1,x_2,y\in[0,1]
\end{equation}

について

\begin{equation}
    \Phi_s(x_1,y)
    =
    \Phi_s(x_2,y)
\end{equation}

が成立するならば、$\Phi_s$ は第1変数に依存しない。

したがって、ある写像

\begin{equation}
    \Psi_s:
    [0,1]
    \longrightarrow
    X
\end{equation}

が存在して、

\begin{equation}
    \Phi_s(x,y)
    =
    \Psi_s(y)
\end{equation}

となる。

この場合、

\begin{equation}
    \mathcal G(s)
    =
    \Psi_s(E_{\mathrm c}(s))
\end{equation}

であり、一般解は離散評価に依存しない。

したがって、$s$ は中間 Solver ではない。

対偶より、中間 Solver ならば

\begin{equation}
    \exists x_1\neq x_2,\ y:
    \qquad
    \Phi_s(x_1,y)
    \neq
    \Phi_s(x_2,y)
\end{equation}

が成立する。

連続評価についても同様に考える。

もし任意の

\begin{equation}
    x,y_1,y_2\in[0,1]
\end{equation}

について

\begin{equation}
    \Phi_s(x,y_1)
    =
    \Phi_s(x,y_2)
\end{equation}

が成立するならば、$\Phi_s$ は第2変数に依存しない。

したがって、ある写像

\begin{equation}
    \Psi'_s:
    [0,1]
    \longrightarrow
    X
\end{equation}

が存在して、

\begin{equation}
    \Phi_s(x,y)
    =
    \Psi'_s(x)
\end{equation}

となる。

この場合、

\begin{equation}
    \mathcal G(s)
    =
    \Psi'_s(E_{\mathrm d}(s))
\end{equation}

であり、一般解は連続評価に依存しない。

したがって、$s$ は中間 Solver ではない。

よって、中間 Solver であるためには、両方の非自明な依存性が
必要である。

逆に、\eqref{eq:H-theorem-discrete} と
\eqref{eq:H-theorem-continuous} がともに成立すれば、
定義 \ref{def:H-middle-solver} の第4条件が成立する。

したがって、

\begin{equation}
    \boxed{
        s\text{ が中間 Solver}
        \iff
        \mathcal G(s)
        \text{ が離散評価・連続評価の双方に
        非自明に依存する}
    }
\end{equation}

が成立する。

\end{proof}


%------------------------------------------------------------
\subsection{(H.17) $\mathcal T(s)$ と二つの評価経路}
\label{subsec:H17}

Solver $s$ に対して、

\begin{equation}
    s=(\mathcal A_s,\mathcal R_s)
\end{equation}

および

\begin{equation}
    \mathcal T(s)
    =
    \mathcal T(\mathcal A_s,\mathcal R_s)
\end{equation}

である。

さらに、

\begin{equation}
    i_s=\mathcal A_s
\end{equation}

であるから、

\begin{equation}
    \boxed{
        \mathcal T(s)
        =
        \mathcal T(i_s,\mathcal R_s)
    }
    \label{eq:H-theory-path}
\end{equation}

と表現できる。

この公理選択 $i_s$ に対して、

\begin{equation}
    D_s=B(i_s)
\end{equation}

が得られる。

したがって、離散評価について

\begin{equation}
    \boxed{
        i_s
        \longrightarrow
        B(i_s)
        \longrightarrow
        E_{\mathrm d}(s)
    }
    \label{eq:H-discrete-path}
\end{equation}

という経路が得られる。

一方、連続評価については

\begin{equation}
    \boxed{
        i_s
        \longrightarrow
        \mu_s
        \longrightarrow
        E_{\mathrm c}(s)
    }
    \label{eq:H-continuous-path}
\end{equation}

という経路が得られる。

したがって、中間 Solver では同一の公理選択 $i_s$ を起点として、

\begin{equation}
    \boxed{
    \begin{aligned}
        i_s
        &\longrightarrow
        D_s
        \longrightarrow
        E_{\mathrm d}(s),
        \\[2mm]
        i_s
        &\longrightarrow
        C_s^{(\mu)}
        \longrightarrow
        E_{\mathrm c}(s)
    \end{aligned}
    }
    \label{eq:H-two-evaluation-paths}
\end{equation}

という二つの評価経路が形成される。


%------------------------------------------------------------
\subsection{(H.18) ECA の離散式と連続式の対応}
\label{subsec:H18}

ECA における離散形式

\begin{equation}
    \mu(i)
    =
    \sum_{b\in B(i)}w(b)
    \label{eq:H-ECA-discrete}
\end{equation}

に

\begin{equation}
    i=i_s
\end{equation}

を代入すると、

\begin{equation}
    \boxed{
        \mu_s^{\mathrm d}(i_s)
        =
        \sum_{b\in B(i_s)}w_s(b)
    }
    \label{eq:H-discrete-form}
\end{equation}

を得る。

一方、ECA における連続形式

\begin{equation}
    \mu(i)
    =
    \int_S
    \mathbf1_i(u)\,d\mu(u)
    \label{eq:H-ECA-continuous}
\end{equation}

に

\begin{equation}
    i=i_s,
    \qquad
    \mu=\mu_s
\end{equation}

を代入すると、

\begin{equation}
    \boxed{
        \mu_s^{\mathrm c}(i_s)
        =
        \int_S
        \mathbf1_{i_s}(u)\,d\mu_s(u)
    }
    \label{eq:H-continuous-form}
\end{equation}

を得る。

したがって、

\begin{equation}
    \boxed{
        E_{\mathrm d}(s)
        =
        \mu_s^{\mathrm d}(i_s)
    }
\end{equation}

および

\begin{equation}
    \boxed{
        E_{\mathrm c}(s)
        =
        \mu_s^{\mathrm c}(i_s)
    }
\end{equation}

である。

この二つの式は、新しい評価原理を導入するものではなく、
同一の公理選択 $i_s$ に対して ECA に既に定義された
離散表現と連続表現を適用したものである。


%------------------------------------------------------------
\subsection{(H.19) 中間 Solver の数学的意味}
\label{subsec:H19}

以上から、中間 Solver は単に

\begin{equation}
    \text{離散集合}\cup\text{連続集合}
\end{equation}

として定義されるものではない。

中間 Solver $s$ においては、

\begin{equation}
    i_s=\mathcal A_s
\end{equation}

を起点として、

\begin{equation}
    D_s=B(i_s)
\end{equation}

という離散評価構造と、

\begin{equation}
    C_s^{(\mu)}=\mu_s
\end{equation}

または

\begin{equation}
    C_s^{(\mathrm{ms})}
    =
    (I_s,\Sigma_{I_s},\mu_s)
\end{equation}

という連続評価構造が構成される。

そして、

\begin{equation}
    E_{\mathrm d}(s)
    =
    \sum_{b\in D_s}w_s(b)
\end{equation}

および

\begin{equation}
    E_{\mathrm c}(s)
    =
    \int_S
    \mathbf1_{i_s}(u)\,d\mu_s(u)
\end{equation}

を介して、

\begin{equation}
    \mathcal G(s)
    =
    \Phi_s
    \left(
        E_{\mathrm d}(s),
        E_{\mathrm c}(s)
    \right)
\end{equation}

が得られる。

従って、中間 Solver の本質は、

\begin{equation}
    \boxed{
    \begin{minipage}{0.85\linewidth}
    同一の Solver に対して離散的評価構造と連続的評価構造が存在し、
    その双方が一般解の決定に非自明に寄与すること。
    \end{minipage}
    }
    \label{eq:H-middle-meaning}
\end{equation}

にある。


%------------------------------------------------------------
\subsection{(H.20) (G) との接続}
\label{subsec:H20}

(G) において一般解は

\begin{equation}
    \mathcal G:
    \mathcal S
    \longrightarrow
    X
\end{equation}

として定義された。

(H) では、この一般解の入力となる Solver の構造を

\begin{equation}
    s
    \longrightarrow
    i_s
    \longrightarrow
    (D_s,C_s)
\end{equation}

という形で具体化した。

したがって、中間 Solver においては、

\begin{equation}
    \boxed{
        s
        \longrightarrow
        i_s
        \longrightarrow
        \left(
            E_{\mathrm d}(s),
            E_{\mathrm c}(s)
        \right)
        \longrightarrow
        \mathcal G(s)
    }
    \label{eq:H-G-connection}
\end{equation}

という構造が得られる。

ここで、$\mathcal G$ 自体は (G) で定義された一般解写像であり、
(H) において新たに定義されるものではない。


%------------------------------------------------------------
\subsection{(H.21) 本節で証明された範囲}
\label{subsec:H21}

本節では、以下を定義または証明した。

\begin{enumerate}[label=\textup{(\arabic*)}]

    \item
    Solver $s$ の公理選択成分

    \begin{equation}
        \mathcal A_s
    \end{equation}

    と ECA の公理選択空間 $I$ を接続するため、

    \begin{equation}
        \mathcal S_I
        =
        \left\{
            s\in\mathcal S
            \mid
            \mathcal A_s\in I
        \right\}
    \end{equation}

    を定義した。

    \item
    ECA に適合する Solver について、

    \begin{equation}
        \mathfrak A(s)
        =
        \mathcal A_s
    \end{equation}

    により公理選択復元写像を構成し、その well-defined 性を
    証明した。

    \item
    したがって、

    \begin{equation}
        i_s
        =
        \mathfrak A(s)
        =
        \mathcal A_s
    \end{equation}

    が一意に定まる。

    \item
    ECA の既存構造から、

    \begin{equation}
        D_s
        =
        B(i_s)
    \end{equation}

    を離散評価構造として定義した。

    \item
    離散評価を

    \begin{equation}
        E_{\mathrm d}(s)
        =
        \sum_{b\in D_s}w_s(b)
    \end{equation}

    と定義し、その有界性を証明した。

    \item
    連続評価について、

    \begin{equation}
        C_s^{(\mu)}
        =
        \mu_s
    \end{equation}

    と

    \begin{equation}
        C_s^{(\mathrm{ms})}
        =
        (I_s,\Sigma_{I_s},\mu_s)
    \end{equation}

    を区別して定義した。

    \item
    離散評価と連続評価を同一の公理選択 $i_s$ に対して
    構成することを示した。

    \item
    一般解が両評価へ非自明に依存することを
    中間 Solver の定義とし、

    \begin{equation}
        s\text{ が中間 Solver}
        \iff
        \mathcal G(s)
        \text{ が離散評価・連続評価の双方に
        非自明に依存する}
    \end{equation}

    を証明した。

\end{enumerate}


%------------------------------------------------------------
\subsection{(H.22) 証明上の前提と ECA 固有の論点}
\label{subsec:H22}

本節の証明によって得られる結果と、
ECA に固有の前提として置かれている事項は区別する必要がある。

特に、

\begin{equation}
    \mathcal A_s\in I
\end{equation}

という Solver--ECA 適合条件は、Solver の一般定義

\begin{equation}
    s=(\mathcal A_s,\mathcal R_s)
\end{equation}

のみから自動的に従うものではない。

したがって、

\begin{equation}
    \boxed{
        \mathcal S_I
        =
        \left\{
            s\in\mathcal S
            \mid
            \mathcal A_s\in I
        \right\}
    }
\end{equation}

を ECA に適合する Solver の部分集合として明示的に扱う。

この区別により、ECA の公理選択空間に属さない Solver や、
必要な理論写像が定義できない Solver を、
中間 Solver の構造へ無理に含めることを避けられる。


%------------------------------------------------------------
\subsection{(H.23) (A)--(G) との統合}
\label{subsec:H23}

(A) において Solver 集合の数学的構造を定義し、

\begin{equation}
    s=(\mathcal A_s,\mathcal R_s)
\end{equation}

という基本表現を与えた。

(B)--(D) において評価空間および一般解

\begin{equation}
    \mathcal G:
    \mathcal S\to X
\end{equation}

を構成した。

(E) において Solver 同型に対する一般解の構造的一意性を扱った。

(F) において離散 Solver の評価を測度的に細分化し、
適切な極限条件のもとで連続評価へ接続した。

(G) において、一般解を

\begin{equation}
    \mathcal S
    \xrightarrow{\mathcal G}
    X
    \xrightarrow{\Psi}
    \mathbb R
\end{equation}

という評価構造として整理した。

そして (H) では、その一般解に入力される Solver の内部構造を

\begin{equation}
    \boxed{
        s
        \longrightarrow
        i_s
        \longrightarrow
        (D_s,C_s)
        \longrightarrow
        (E_{\mathrm d}(s),E_{\mathrm c}(s))
        \longrightarrow
        \mathcal G(s)
    }
    \label{eq:H-complete-chain}
\end{equation}

として具体化した。

特に中間 Solver では、

\begin{equation}
    \boxed{
        \mathcal G(s)
        =
        \Phi_s
        \left(
            E_{\mathrm d}(s),
            E_{\mathrm c}(s)
        \right)
    }
\end{equation}

において双方の評価が非自明に寄与する。


%------------------------------------------------------------
\subsection{(H.24) 本節の結論}
\label{subsec:H24}

以上により、(H) における中間 Solver の構造は、

\begin{equation}
    \boxed{
        s=(\mathcal A_s,\mathcal R_s)
    }
\end{equation}

から始まり、

\begin{equation}
    \boxed{
        i_s
        =
        \mathfrak A(s)
        =
        \mathcal A_s
    }
\end{equation}

を得て、

\begin{equation}
    \boxed{
        D_s=B(i_s)
    }
\end{equation}

および

\begin{equation}
    \boxed{
        C_s^{(\mu)}=\mu_s
        \qquad
        \text{または}
        \qquad
        C_s^{(\mathrm{ms})}
        =
        (I_s,\Sigma_{I_s},\mu_s)
    }
\end{equation}

を構成する。

さらに、

\begin{equation}
    \boxed{
        E_{\mathrm d}(s)
        =
        \sum_{b\in D_s}w_s(b)
    }
\end{equation}

および

\begin{equation}
    \boxed{
        E_{\mathrm c}(s)
        =
        \int_S
        \mathbf1_{i_s}(u)\,d\mu_s(u)
    }
\end{equation}

を得て、

\begin{equation}
    \boxed{
        \mathcal G(s)
        =
        \Phi_s
        \left(
            E_{\mathrm d}(s),
            E_{\mathrm c}(s)
        \right)
    }
\end{equation}

として一般解へ接続する。

したがって、

\begin{equation}
    \boxed{
        \text{Solver}
        \longrightarrow
        \text{公理選択}
        \longrightarrow
        \text{離散・連続評価構造}
        \longrightarrow
        \text{一般解}
    }
\end{equation}

という構造が得られる。

このとき中間 Solver とは、

\begin{equation}
    \boxed{
    \begin{minipage}{0.85\linewidth}
    同一の Solver 内に離散的評価構造と連続的評価構造が存在し、
    その双方が一般解の決定に非自明に寄与する Solver
    \end{minipage}
    }
\end{equation}

として厳密に定義される。



\end{document}
